% ======================================================================
% sn-article.tex  --  Mathematical Geosciences submission
%   Springer Nature sn-jnl class  (sn-mathphys-ay = author-year refs)
% Title: A Physics-Informed Deep Learning Framework for Solving the
%        2D Shallow Water Equations in Non-Stationary Flood Environments
% Author: Kouao Laurent Kouadio  (ORCID: 0000-0001-7259-7254)
% ======================================================================

%% Math and Physical Sciences Author-Year Reference Style
\documentclass[pdflatex,sn-mathphys-ay]{sn-jnl}

%%---standard packages--------------------------------------------------
\usepackage{graphicx}
\usepackage{amsmath,amssymb,amsfonts}
\usepackage{mathrsfs}            % \mathscr
\usepackage{booktabs}            % \toprule / \midrule / \bottomrule
\usepackage{multirow}
\usepackage[title]{appendix}    % \begin{appendices}
\usepackage{xcolor}              % \textcolor in figure captions
\usepackage{enumerate}
\usepackage{textcomp}
\usepackage{url}                 % robust URLs in declarations
\usepackage{float}               % float support; floats placed at end
\graphicspath{{figures/}{./}}
%%----------------------------------------------------------------------

%%---theorem environments (sn-jnl provides thmstyleone/two/three)------
\theoremstyle{thmstyleone}
\newtheorem{theorem}{Theorem}
\newtheorem{lemma}[theorem]{Lemma}
\newtheorem{corollary}[theorem]{Corollary}
\newtheorem{proposition}[theorem]{Proposition}

\theoremstyle{thmstylethree}
\newtheorem{definition}{Definition}

\theoremstyle{thmstyletwo}
\newtheorem{remark}{Remark}
%%----------------------------------------------------------------------

\raggedbottom

% =======================================================================
\begin{document}

\title[Physics-Informed Reservoir Learning for Flood Modelling]{%
  Physics-Informed Reservoir Learning for Shallow-Water Flood Modelling}

\author*[1,2]{\fnm{Kouao Laurent} \sur{Kouadio}}%
 %% \orcid{0000-0001-7259-7254}}\email{etanoyau@gmail.com}

\affil*[1]{\orgdiv{School of Geosciences and Info-Physics},
  \orgname{Central South University},
  \orgaddress{\city{Changsha}, \state{Hunan 410083}, \country{China}}}

\affil[2]{\orgdiv{Hunan Key Laboratory of Nonferrous Resources and
  Geological Hazards Exploration},
  \orgname{Central South University},
  \orgaddress{\city{Changsha}, \state{Hunan 410083}, \country{China}}}


\abstract{Flood prediction in data-sparse regions requires models that are both computationally efficient and physically consistent. Classical
two-dimensional Shallow Water Equation (SWE) solvers preserve mass and
momentum conservation but remain costly for large-scale or ensemble
forecasting, whereas purely data-driven surrogates may violate
hydraulic constraints under rare or non-stationary flood regimes. We
introduce the Physically-Aware Deep Reservoir Network
(\textsc{PADR-Net}), a physics-informed reservoir-learning framework
that couples a contractive Echo State reservoir with a shallow-water
residual penalty. The reservoir encodes hydroclimatic memory, while the
physics term constrains predicted flood states toward the admissible
SWE solution manifold. We establish three properties of
the framework: strict hyperbolicity of the SWE system for positive
water depth, fading memory of the reservoir state under the
contractivity condition, and a residual-distance bound showing that the
distance to the SWE solution manifold decreases as
$O(\lambda^{-1/2})$ with increasing physics weight $\lambda$. The
framework is evaluated on controlled synthetic flood regimes and on a
harmonized archive of 243 African flood events from 2000 to 2024
across West Africa, East Africa, and Southern Africa. In the synthetic
benchmark, \textsc{PADR-Net} reproduced monsoon, flash/cyclone, and
seasonal-memory flood responses with mean CSI = 0.86 and mean TSS =
0.92. Antecedent memory descriptors improved flood-severity ranking from
$\rho_s=0.19$ to $\rho_s=0.61$, while exposure and hydrodynamic
descriptors increased high-impact event discrimination from PR-AUC =
0.403 to 0.703 and improved flood-depth reconstruction from
NSE$_\mathrm{depth}=0.38$ to $0.64$. Transfer experiments showed that depth response
generalizes more consistently than impact ranking across regions and
years. These results show that flood burden is not a function of
rainfall intensity alone, but emerges from the coupling between
meteorological forcing, antecedent catchment state, terrain-mediated
routing, and ex-ante human exposure.}

\keywords{Physics-informed learning,
          Shallow Water Equations,
          Reservoir computing,
          Flood inundation modeling,
          Flood impact prediction,
          Hydroclimatic memory,
          African flood events}



%%\pacs[MSC Classification]{86A05, 68T07, 65M99}

\maketitle

% =======================================================================
\section{Introduction}
\label{sec:intro}

Floods remain among the most damaging and spatially pervasive natural
hazards, affecting populations, infrastructure, agricultural systems,
and regional economies across a wide range of climatic and
geomorphological settings.  Their impacts are increasingly shaped by
the combined effects of climatic non-stationarity, land-use change,
rapid urban expansion, and the continued settlement of flood-prone
terrain.  Satellite observations have shown that human exposure to
flooding is not merely a passive consequence of hazard occurrence but
is actively changing through the redistribution of population and built
assets into areas where inundation has been observed
\citep{Jongman2012,Ceola2014,Tellman2021,Rentschler2022}.  At the same time, global projections indicate
that river flood risk is sensitive to both climate forcing and
socio-economic development pathways, with future changes varying
substantially across regions and hydrological regimes
\citep{Hirabayashi2013,Winsemius2016,Alfieri2016}.  These interacting
drivers make flood prediction and impact estimation difficult,
particularly in data-scarce regions where long hydrometric records,
high-resolution topography, and reliable impact observations are often
limited.

A persistent limitation in conventional flood assessment is the
assumption that rainfall extremes can be translated directly into flood
hazard or reported human impact.  Although precipitation provides the
primary atmospheric forcing for many flood events, its relationship
with inundation is strongly mediated by antecedent soil moisture,
catchment storage, river-network routing, floodplain connectivity,
urban drainage, reservoir regulation, and terrain structure.  This
nonlinear disconnect between precipitation anomalies and regional flood
activity has been described through the concept of ``floodiness'',
which emphasizes that flood occurrence and magnitude cannot be
reliably inferred from rainfall statistics alone
\citep{Stephens2015}.  The problem becomes even more complex in
mountainous and high-energy environments, where flood disasters may
arise from cascading processes such as glacial lake outburst floods,
landslides, sediment entrainment, or dam-break-like releases rather
than from local rainfall intensity alone \citep{Sattar2025}.  A
general flood-impact model must therefore represent not only the
meteorological trigger but also the hydrological memory and physical
routing mechanisms that transform rainfall or water release into
spatially distributed inundation.

Physically based hydrodynamic models provide the most direct route for
representing these processes.  Models based on the two-dimensional
shallow-water equations enforce conservation of mass and momentum and
can reproduce flow acceleration, wetting and drying, floodplain
storage, and terrain-controlled propagation.  Efficient inertial and
local-inertial formulations have made large-scale flood simulation more
tractable while retaining the essential structure of shallow-water
dynamics \citep{Bates2010,Almeida2013,Xia2017,Bates2022}.  These developments
have supported continental and global flood hazard modelling, including
high-resolution inundation products and multi-peril frameworks
designed to represent fluvial, pluvial, and coastal flooding
\citep{Sampson2015,Dottori2016,Trigg2016,Wing2019,Bates2023,Wing2024}.  Operational and
pre-operational forecasting systems have also demonstrated the value of
large-scale hydrometeorological modelling for early warning
\citep{Thielen2009,Alfieri2013,Emerton2016,Alfieri2019,Pappenberger2015,Najafi2024ImpactBasedFloodWarning}.  Nevertheless,
high-resolution hydrodynamic simulation remains computationally
expensive when repeated over many events, ensemble members, climate
scenarios, or uncertainty realizations.  This limitation is especially
restrictive for real-time prediction and for continental-scale impact
analysis in regions where local hydraulic calibration data are sparse.

Machine-learning models offer an attractive complement because they can
learn nonlinear relationships from high-dimensional hydroclimatic,
topographic, and exposure data with far lower computational cost after
training.  Recurrent architectures, particularly long short-term memory
networks, have substantially improved rainfall--runoff prediction and
large-sample hydrological modelling \citep{Hochreiter1997,Kratzert2019,Kratzert2023}.
Recent work has further shown that deep learning can generalize useful
hydrological information across ungauged or poorly gauged basins when
large and diverse training data are available \citep{Nearing2024}.
However, purely data-driven models remain vulnerable when confronted
with non-stationary extremes, unseen hydrological regimes, or sparse
training examples.  Without physical constraints, they may reproduce
historical correlations while violating basic hydraulic principles,
such as water conservation, terrain-consistent routing, or momentum
balance.  This weakness is particularly problematic for flood-impact
estimation, where the most important events are often rare,
out-of-distribution, and controlled by coupled hazard--exposure
mechanisms.

The challenge is therefore not to choose between physical modelling and
machine learning, but to construct a framework in which the two forms
of information reinforce one another.  A physically guided surrogate
should retain the speed and flexibility of learning-based prediction
while remaining anchored to the governing equations that define
admissible flood states.  In this study, we develop such a framework
through a Physically-Aware Deep Reservoir Network, denoted
\textsc{PADR-Net}.  The formulation follows the broader physics-informed
learning principle of embedding governing-equation residuals in the
training objective \citep{RAISSI2019686}.  The model uses a reservoir
architecture to capture
nonlinear hydroclimatic memory from sequences of rainfall, antecedent
wetness, river-state indicators, terrain descriptors, and exposure
covariates.  Its predictions are then regularized by residuals derived
from the shallow-water balance, so that the learned trajectories are
encouraged to remain close to the physically admissible solution
manifold.  This formulation is designed to address the central
difficulty of flood-impact prediction: damaging floods emerge from the
interaction between meteorological forcing, catchment memory,
hydrodynamic propagation, and the spatial concentration of people and
assets in flood-prone environments.

The development of the framework proceeds by linking the
structure of the shallow-water equations to the stability of the
learning model.  The governing system is shown to be strictly
hyperbolic for positive water depth, which establishes the physical
basis for representing propagating flood waves through real-valued
characteristic speeds.  The reservoir state map is then formulated as a
contractive dynamical system under the echo-state condition
\citep{Jaeger2001EchoState,Herbert2004,LUKOSEVICIUS2009127}, ensuring
that the learned state is controlled by the forcing history rather than
by arbitrary initialization.  Finally, the physics-informed objective is
shown to impose a quantitative bound on the distance between the
predicted trajectory and the shallow-water solution manifold, with the
bound decreasing as the physics weight increases.  These results give a
formal interpretation to the proposed regularization: the physical loss
is not only a numerical penalty but a mechanism that constrains the
surrogate model toward mass- and momentum-consistent flood dynamics.

The empirical design follows the same logic.  We first evaluate
\textsc{PADR-Net} under a controlled synthetic flood-event benchmark in
which rainfall forcing, antecedent wetness, terrain routing, exposure,
and impact generation are known.  This benchmark allows us to test
whether the model recovers the intended hierarchy among rainfall,
memory, hydrodynamics, and exposure.  We then apply the framework to a
harmonized archive of 243 African flood events from 2000--2024,
covering contrasting flood-prone regions including the Niger--Benue
system, the Nile headwaters and Sudan--Ethiopia corridor, and the
Limpopo--Zambezi domain.  Africa provides a particularly important
test bed because flood prediction is constrained by heterogeneous
observation networks, strong hydroclimatic seasonality, rapidly
changing exposure, and large regional differences in rainfall regimes
and basin response
\citep{DiBaldassarre2010,Thiemig2015,Nikulin2012,Dosio2017,Bodian2018,Dembele2020}.
The experiments are designed to test not only whether the model fits
past events, but whether it transfers across regions and years under
conditions where rainfall-only predictors are expected to fail.

By integrating hydrodynamic constraints, reservoir-based memory, and
ex-ante exposure features, this work advances a flood-impact modelling
framework that is both physically interpretable and computationally
efficient.  The central hypothesis is that robust prediction of flood
burden requires moving beyond precipitation intensity as the dominant
explanatory variable.  Instead, flood impacts must be understood as the
outcome of coupled climate forcing, antecedent catchment state,
terrain-mediated water routing, and human exposure.  The following
sections formalize this framework, describe the synthetic and African
experiments, and evaluate whether physics-aware learning improves
spatial inundation prediction and impact ranking under transfer across
diverse flood regimes; the detailed data provenance, harmonisation
rules, leakage safeguards, and validation partitions are reported in
Appendix~\ref{app:data_experimental_details}.

% =======================================================================
\section{Methodology}
\label{sec:methodology}

The methodology is designed to test whether flood dynamics and
reported impacts can be explained by event-scale rainfall alone, or
whether reliable prediction requires an explicit coupling between
hydrodynamic physics, antecedent hydroclimatic memory, terrain
structure, and pre-event exposure.  Each flood is represented as an
event--region record indexed by $i=1,\ldots,N$, associated with a
regional domain $\Omega_{r_i}\subset\mathbb{R}^{2}$, an event year
$y_i$, and an event time window
\begin{equation}
  \mathcal{T}_{i}
  =
  \{t:s_i\leq t\leq e_i\},
  \label{eq:event_window}
\end{equation}
where $s_i$ and $e_i$ denote the start and end dates of the event.
Events with incomplete dates are retained for descriptive summaries,
but they are excluded from event-window hydroclimatic modelling unless
a physically defensible temporal window can be assigned.  The resulting
event catalogue combines satellite-observed flood information, disaster
impact records, reanalysis hydroclimatic fields, terrain descriptors,
and gridded population and built-up layers
\citep{EMDAT2024,Tellman2021,Stevens2015,Pesaresi2023,Hersbach2020}.  Satellite flood maps are
used as observational targets and validation data, whereas exposure
features used for prediction are computed only from information
available before the event, so that the realized inundation polygon does
not enter the predictor set.  The corresponding data provenance,
event-window rules, and harmonisation safeguards are detailed in
Appendices~\ref{app:data_sources}--\ref{app:event_harmonisation}.

For each record, the forcing field is constructed from precipitation,
antecedent wetness, temperature, and terrain variables sampled over
$\Omega_{r_i}$.  The antecedent precipitation accumulated over a memory
scale $\tau$ is written as
\begin{equation}
  P^{\mathrm{ant}}_{i,\tau}
  =
  \frac{
  \displaystyle
  \sum_{t=s_i-\tau}^{s_i-1}
  \sum_{g\in\Omega_{r_i}} A_g P_{t,g}
  }{
  \displaystyle
  \sum_{g\in\Omega_{r_i}} A_g
  },
  \qquad
  \tau\in\{7,14,30,60\},
  \label{eq:antecedent_precip}
\end{equation}
where $A_g$ is the area of grid cell $g$ and $P_{t,g}$ is daily
precipitation.  Soil-moisture memory is computed analogously, replacing
$P_{t,g}$ by the corresponding soil-moisture variable.  Where observed
or routed discharge is available, the pre-event river state is
standardized relative to the seasonal climatology,
\begin{equation}
  R^{\mathrm{ant}}_{i,\tau}
  =
  \frac{1}{\tau}
  \sum_{t=s_i-\tau}^{s_i-1}
  \frac{
    Q_{t,h(i)}-\mu_{h(i),m(t)}
  }{
    \sigma_{h(i),m(t)}+\varepsilon
  },
  \label{eq:river_state}
\end{equation}
where $Q_{t,h(i)}$ is discharge at the hydrologically relevant gauge or
routing node $h(i)$, $\mu_{h(i),m(t)}$ and $\sigma_{h(i),m(t)}$ are the
mean and standard deviation for calendar month $m(t)$, and
$\varepsilon>0$ prevents division by zero.  These variables make the
state of the catchment before the event explicit, which is essential
because regional floodiness is often controlled by storage, routing,
and saturation effects rather than by rainfall magnitude alone
\citep{Seneviratne2010,Stephens2015}.

The physical component of the framework is based on the two-dimensional
shallow-water equations, written in conservative form over
$\Omega_{r_i}$ as
\begin{equation}
  \frac{\partial \mathbf{q}}{\partial t}
  +
  \frac{\partial \mathbf{f}}{\partial x}
  +
  \frac{\partial \mathbf{g}}{\partial y}
  =
  \mathbf{R}
  +
  \mathbf{S}_{b}
  +
  \mathbf{S}_{f}.
  \label{eq:swe_main}
\end{equation}
Here $\mathbf{q}$ is the vector of conserved variables, while
$\mathbf{f}$ and $\mathbf{g}$ are the fluxes in the $x$ and $y$
directions:
\begin{equation}
  \mathbf{q}
  =
  \begin{bmatrix}
    h \\ uh \\ vh
  \end{bmatrix},
  \quad
  \mathbf{f}
  =
  \begin{bmatrix}
    uh \\
    u^{2}h+\tfrac{1}{2}g_{0}h^{2} \\
    uvh
  \end{bmatrix},
  \quad
  \mathbf{g}
  =
  \begin{bmatrix}
    vh \\
    uvh \\
    v^{2}h+\tfrac{1}{2}g_{0}h^{2}
  \end{bmatrix},
  \label{eq:swe_fluxes}
\end{equation}
where $h>0$ is water depth, $u$ and $v$ are depth-averaged velocities,
and $g_{0}$ is gravitational acceleration
\citep{Toro2001,Bates2010,Liang2010,Bates2022}.
The source terms represent mass exchange, bed slope, and friction:
\begin{equation}
  \mathbf{R}
  =
  \begin{bmatrix}
    P-I-D+Q^{\mathrm{lat}} \\ 0 \\ 0
  \end{bmatrix},
  \quad
  \mathbf{S}_{b}
  =
  \begin{bmatrix}
    0 \\
    -g_{0}h\,\partial b/\partial x \\
    -g_{0}h\,\partial b/\partial y
  \end{bmatrix},
  \quad
  \mathbf{S}_{f}
  =
  \begin{bmatrix}
    0 \\
    -\tau_{bx}/\rho \\
    -\tau_{by}/\rho
  \end{bmatrix}.
  \label{eq:swe_sources}
\end{equation}
In this formulation, $P$ is precipitation, $I$ is infiltration, $D$ is
drainage loss, $Q^{\mathrm{lat}}$ denotes any lateral inflow, $b$ is
bed elevation, and $\rho$ is water density.  The friction stresses are
evaluated using the Manning closure,
\begin{equation}
  \tau_{bx}
  =
  \rho C_f u\sqrt{u^2+v^2},
  \qquad
  \tau_{by}
  =
  \rho C_f v\sqrt{u^2+v^2},
  \qquad
  C_f
  =
  \frac{g_{0}n^2}{h^{1/3}},
  \label{eq:manning}
\end{equation}
where $n$ is the spatially distributed Manning roughness coefficient
\citep{Akan2006,Arcement1989,Almeida2013,Xia2017}.  The bed-elevation field $b(x,y)$ is
derived from high-resolution terrain data, with positive elevation
biases from vegetation and built structures corrected where possible
\citep{Farr2007,Yamazaki2017,Hawker2022}.  The equations are written in full
conservative form because this provides a common physical residual for
training; when simplified hydrodynamic products are used for
benchmarking, they are treated as reference simulations rather than as
the governing definition of the physics loss.

The hyperbolic structure of the shallow-water flux provides the basis for physically admissible flood-wave propagation
\citep{Toro2001,Roe1981}.  This property is stated formally in
Theorem~\ref{thm:swe_hyperbolicity}.

\begin{theorem}[Strict hyperbolicity of the shallow-water system]
\label{thm:swe_hyperbolicity}
For positive water depth $h>0$, the conservative shallow-water system
in Eq.~\eqref{eq:swe_main} is strictly hyperbolic in each coordinate
direction.  Consequently, the system admits real characteristic wave
speeds and supports physically meaningful flood-wave propagation.
\end{theorem}

\begin{proof}
Writing $\mathbf{q}=(h,q_x,q_y)^\top$, with $q_x=uh$ and $q_y=vh$,
gives $u=q_x/h$ and $v=q_y/h$.  Differentiating the $x$-flux with
respect to $\mathbf{q}$ yields
\begin{equation}
  \mathbf{A}
  =
  \frac{\partial\mathbf{f}}{\partial\mathbf{q}}
  =
  \begin{bmatrix}
    0 & 1 & 0 \\
    c^{2}-u^{2} & 2u & 0 \\
    -uv & v & u
  \end{bmatrix},
  \qquad
  c=\sqrt{g_{0}h}.
  \label{eq:jacobian}
\end{equation}
The characteristic polynomial factorizes as
$(\mu-u)\{(\mu-u)^2-c^2\}=0$, so the three wave speeds are
\begin{equation}
  \lambda_1=u-c,
  \qquad
  \lambda_2=u,
  \qquad
  \lambda_3=u+c.
  \label{eq:eigenvalues}
\end{equation}
Since $h>0$ implies $c>0$, the eigenvalues are real and distinct.
The same argument applies in the $y$ direction, with $v$ replacing
$u$.  Thus, the governing system is strictly hyperbolic wherever water
depth is positive.
\end{proof}

This property explains why the formulation can represent surge fronts
and rapidly propagating flood waves, while also providing a physically
meaningful residual against which the learning model can be
constrained.

%%%%%%%%%%%%%%%%%%%%%%%%==============================================

To approximate the evolving flood state while retaining the physical
constraints imposed by Eq.~\eqref{eq:swe_main}, we introduce a
physics-informed deep reservoir formulation, denoted \textsc{PADR-Net}
\citep{RAISSI2019686,Jaeger2001EchoState,LUKOSEVICIUS2009127}.
The model is designed to map a sequence of hydroclimatic and
geomorphic forcings to a physically admissible flood state.  At each
time step $t$, the input vector
$\mathbf{u}_{i,t}\in\mathbb{R}^{N_{\mathrm{in}}}$ contains event
precipitation, antecedent wetness, river-state indicators, terrain
gradient, flow-accumulation descriptors, seasonal variables, and
regional hydroclimatic indicators.  These variables drive a
high-dimensional reservoir state
$\mathbf{x}_{i,t}\in\mathbb{R}^{N_{\mathrm{res}}}$ according to
\begin{equation}
  \mathbf{x}_{i,t}
  =
  \sigma
  \left(
    \mathbf{W}_{\mathrm{in}}\mathbf{u}_{i,t}
    +
    \mathbf{W}_{\mathrm{res}}\mathbf{x}_{i,t-1}
    +
    \mathbf{b}_{\mathrm{res}}
  \right),
  \label{eq:padr_state}
\end{equation}
where $\mathbf{W}_{\mathrm{in}}$ and $\mathbf{W}_{\mathrm{res}}$ are
fixed random matrices, $\mathbf{b}_{\mathrm{res}}$ is a reservoir bias
term, and $\sigma(\cdot)$ is a bounded nonlinear activation function.
Only the output mapping is optimized.  The predicted conservative
hydrodynamic state is
\begin{equation}
  \hat{\mathbf{q}}_{i,t,g}
  =
  \begin{bmatrix}
    \hat{h}_{i,t,g} \\
    \widehat{uh}_{i,t,g} \\
    \widehat{vh}_{i,t,g}
  \end{bmatrix}
  =
  \mathbf{W}_{\mathrm{out}}
  \mathbf{x}_{i,t,g},
  \label{eq:padr_output}
\end{equation}
where $g$ denotes a grid cell, $\hat{h}_{i,t,g}$ is predicted water
depth, and $\widehat{uh}_{i,t,g}$ and $\widehat{vh}_{i,t,g}$ are the
predicted momentum components.  The corresponding velocity components
are recovered as
\begin{equation}
  \hat{u}_{i,t,g}
  =
  \frac{\widehat{uh}_{i,t,g}}
  {\hat{h}_{i,t,g}+\varepsilon_h},
  \qquad
  \hat{v}_{i,t,g}
  =
  \frac{\widehat{vh}_{i,t,g}}
  {\hat{h}_{i,t,g}+\varepsilon_h},
  \label{eq:padr_velocity}
\end{equation}
with $\varepsilon_h>0$ introduced to avoid numerical instability in
nearly dry cells.

The usefulness of the reservoir representation depends on whether the
state is governed by the input history rather than by arbitrary initial
conditions.  This fading-memory property is formalized in
Lemma~\ref{lem:echo_state}.

\begin{lemma}[Contractive reservoir and fading memory]
\label{lem:echo_state}
Consider the reservoir update in Eq.~\eqref{eq:padr_state}, and assume
that the activation function is Lipschitz continuous with constant
$L_{\sigma}$.  If
$L_{\sigma}\|\mathbf{W}_{\mathrm{res}}\|_{2}<1$, then the reservoir
map is contractive.  Consequently, for two states driven by the same
input sequence, the influence of arbitrary initial conditions decays
geometrically.
\end{lemma}

\begin{proof}
This property is obtained by enforcing a contractive reservoir.  If
$L_{\sigma}$ is the Lipschitz constant of the activation function and
\begin{equation}
  \alpha_{\mathrm{res}}
  =
  L_{\sigma}
  \left\|
    \mathbf{W}_{\mathrm{res}}
  \right\|_{2}
  <1,
  \label{eq:reservoir_contractivity_condition}
\end{equation}
then, for two states $\mathbf{x}_{i,t}$ and $\mathbf{x}'_{i,t}$ driven
by the same input sequence, Eq.~\eqref{eq:padr_state} gives
\begin{equation}
  \left\|
    \mathbf{x}_{i,t}
    -
    \mathbf{x}'_{i,t}
  \right\|_{2}
  \leq
  \alpha_{\mathrm{res}}
  \left\|
    \mathbf{x}_{i,t-1}
    -
    \mathbf{x}'_{i,t-1}
  \right\|_{2}.
  \label{eq:reservoir_one_step}
\end{equation}
Repeated application over $t$ time steps yields
\begin{equation}
  \left\|
    \mathbf{x}_{i,t}
    -
    \mathbf{x}'_{i,t}
  \right\|_{2}
  \leq
  \alpha_{\mathrm{res}}^{t}
  \left\|
    \mathbf{x}_{i,0}
    -
    \mathbf{x}'_{i,0}
  \right\|_{2},
  \label{eq:reservoir_fading_memory}
\end{equation}
which converges to zero as $t$ increases.  Hence, after a sufficiently
long wash-in period, the reservoir state is determined by the forcing
history.
\end{proof}

This fading-memory property is important for flood modelling because
flood response is controlled not only by instantaneous rainfall but
also by accumulated catchment wetness, river routing, and storage
conditions inherited from earlier times
\citep{Jaeger2001EchoState,Herbert2004}.

The predicted state in Eq.~\eqref{eq:padr_output} is constrained by
the shallow-water balance through a residual operator.  Substitution of
$\hat{\mathbf{q}}_{i,t,g}$ into
Eqs.~\eqref{eq:swe_main}--\eqref{eq:swe_sources} gives
\begin{equation}
  \mathcal{F}
  \left(
    \hat{\mathbf{q}}_{i,t,g}
  \right)
  =
  \begin{bmatrix}
    \mathcal{F}^{h}_{i,t,g} \\
    \mathcal{F}^{x}_{i,t,g} \\
    \mathcal{F}^{y}_{i,t,g}
  \end{bmatrix},
  \label{eq:padr_residual_vector}
\end{equation}
where the continuity residual is
\begin{equation}
  \mathcal{F}^{h}_{i,t,g}
  =
  \frac{\partial \hat{h}}{\partial t}
  +
  \frac{\partial (\hat{u}\hat{h})}{\partial x}
  +
  \frac{\partial (\hat{v}\hat{h})}{\partial y}
  -
  \left(
    P-I-D+Q^{\mathrm{lat}}
  \right),
  \label{eq:padr_mass_residual}
\end{equation}
and the two momentum residuals are
\begin{align}
  \mathcal{F}^{x}_{i,t,g}
  &=
  \frac{\partial (\hat{u}\hat{h})}{\partial t}
  +
  \frac{\partial}{\partial x}
  \left(
    \hat{u}^{2}\hat{h}
    +
    \frac{1}{2}g_{0}\hat{h}^{2}
  \right)
  +
  \frac{\partial}{\partial y}
  \left(
    \hat{u}\hat{v}\hat{h}
  \right)
  +
  g_{0}\hat{h}
  \frac{\partial b}{\partial x}
  +
  \frac{\hat{\tau}_{bx}}{\rho},
  \label{eq:padr_xmomentum_residual}
  \\
  \mathcal{F}^{y}_{i,t,g}
  &=
  \frac{\partial (\hat{v}\hat{h})}{\partial t}
  +
  \frac{\partial}{\partial x}
  \left(
    \hat{u}\hat{v}\hat{h}
  \right)
  +
  \frac{\partial}{\partial y}
  \left(
    \hat{v}^{2}\hat{h}
    +
    \frac{1}{2}g_{0}\hat{h}^{2}
  \right)
  +
  g_{0}\hat{h}
  \frac{\partial b}{\partial y}
  +
  \frac{\hat{\tau}_{by}}{\rho}.
  \label{eq:padr_ymomentum_residual}
\end{align}
The predicted friction stresses are computed as
\begin{equation}
  \hat{\tau}_{bx}
  =
  \rho \hat{C}_{f}
  \hat{u}
  \sqrt{\hat{u}^{2}+\hat{v}^{2}},
  \qquad
  \hat{\tau}_{by}
  =
  \rho \hat{C}_{f}
  \hat{v}
  \sqrt{\hat{u}^{2}+\hat{v}^{2}},
  \qquad
  \hat{C}_{f}
  =
  \frac{g_{0}n^{2}}
  {(\hat{h}+\varepsilon_h)^{1/3}}.
  \label{eq:padr_friction}
\end{equation}
In practice, the derivatives in
Eqs.~\eqref{eq:padr_mass_residual}--\eqref{eq:padr_ymomentum_residual}
are evaluated at collocation points using differentiable numerical
operators.  For a generic variable $\phi$, the discrete temporal and
spatial derivatives are written as
\begin{equation}
  \partial_{t}\phi_{i,t,g}
  \approx
  \frac{
    \phi_{i,t+\Delta t,g}
    -
    \phi_{i,t,g}
  }{\Delta t},
  \qquad
  \partial_{x}\phi_{i,t,g}
  \approx
  \frac{
    \phi_{i,t,g^{+}_{x}}
    -
    \phi_{i,t,g^{-}_{x}}
  }{2\Delta x},
  \qquad
  \partial_{y}\phi_{i,t,g}
  \approx
  \frac{
    \phi_{i,t,g^{+}_{y}}
    -
    \phi_{i,t,g^{-}_{y}}
  }{2\Delta y}.
  \label{eq:padr_discrete_derivatives}
\end{equation}
This makes the physics term computable on the same grid used for the
satellite and hydrodynamic validation products.

The learning problem can be written as a regularized
pseudo-posterior.  Let
$\mathcal{D}_{\mathrm{obs}}$ be the set of observed or benchmark flood
states and $\mathcal{D}_{\mathrm{col}}$ the set of collocation points
where the shallow-water residual is evaluated.  The parameter block is
defined as
\begin{equation}
  \Theta
  =
  \left(
    \mathbf{W}_{\mathrm{out}},
    \sigma^{2}_{y},
    \sigma^{2}_{F},
    \eta
  \right),
  \label{eq:padr_parameter_block}
\end{equation}
where $\sigma^{2}_{y}$ is the observational-error variance,
$\sigma^{2}_{F}$ is the physics-residual variance, and $\eta$ controls
the output-weight regularization.  The data model is
\begin{equation}
  \mathbf{y}^{\mathrm{obs}}_{i,t,g}
  =
  \hat{\mathbf{q}}_{i,t,g}
  +
  \boldsymbol{\epsilon}_{i,t,g},
  \qquad
  \boldsymbol{\epsilon}_{i,t,g}
  \sim
  \mathcal{N}
  \left(
    \mathbf{0},
    \sigma^{2}_{y}\mathbf{I}
  \right),
  \label{eq:padr_data_model}
\end{equation}
whereas the physical residual is treated as a zero-centred
pseudo-observation:
\begin{equation}
  \mathbf{0}
  =
  \mathcal{F}
  \left(
    \hat{\mathbf{q}}_{i,t,g}
  \right)
  +
  \boldsymbol{\xi}_{i,t,g},
  \qquad
  \boldsymbol{\xi}_{i,t,g}
  \sim
  \mathcal{N}
  \left(
    \mathbf{0},
    \sigma^{2}_{F}\mathbf{I}
  \right).
  \label{eq:padr_physics_model}
\end{equation}
Combining Eqs.~\eqref{eq:padr_data_model} and
\eqref{eq:padr_physics_model} with a Gaussian prior on
$\mathbf{W}_{\mathrm{out}}$ gives the joint density
\begin{align}
  \pi
  \left(
    \Theta
    \mid
    \mathcal{D}_{\mathrm{obs}},
    \mathcal{D}_{\mathrm{col}}
  \right)
  &\propto
  \left[
    \prod_{(i,t,g)\in\mathcal{D}_{\mathrm{obs}}}
    \mathcal{N}
    \left(
      \mathbf{y}^{\mathrm{obs}}_{i,t,g}
      \mid
      \mathbf{W}_{\mathrm{out}}\mathbf{x}_{i,t,g},
      \sigma^{2}_{y}\mathbf{I}
    \right)
  \right]
  \nonumber \\
  &\quad\times
  \left[
    \prod_{(i,t,g)\in\mathcal{D}_{\mathrm{col}}}
    \mathcal{N}
    \left(
      \mathbf{0}
      \mid
      \mathcal{F}
      \left(
        \mathbf{W}_{\mathrm{out}}\mathbf{x}_{i,t,g}
      \right),
      \sigma^{2}_{F}\mathbf{I}
    \right)
  \right]
  \nonumber \\
  &\quad\times
  \mathcal{N}
  \left(
    \mathrm{vec}
    \left(
      \mathbf{W}_{\mathrm{out}}
    \right)
    \mid
    \mathbf{0},
    \eta^{-1}\mathbf{I}
  \right)
  p
  \left(
    \sigma^{2}_{y}
  \right)
  p
  \left(
    \sigma^{2}_{F}
  \right)
  p(\eta).
  \label{eq:padr_joint_density}
\end{align}
This expression makes explicit the three sources of information used
by the model: observed flood states, hydrodynamic conservation laws,
and parameter regularity.

Taking the negative logarithm of
Eq.~\eqref{eq:padr_joint_density}, discarding constants independent of
$\mathbf{W}_{\mathrm{out}}$, and collecting terms gives
\begin{align}
  \mathcal{J}
  \left(
    \mathbf{W}_{\mathrm{out}}
  \right)
  &=
  \frac{1}{2\sigma^{2}_{y}}
  \sum_{(i,t,g)\in\mathcal{D}_{\mathrm{obs}}}
  \left\|
    \mathbf{y}^{\mathrm{obs}}_{i,t,g}
    -
    \mathbf{W}_{\mathrm{out}}\mathbf{x}_{i,t,g}
  \right\|_{2}^{2}
  \nonumber \\
  &\quad+
  \frac{1}{2\sigma^{2}_{F}}
  \sum_{(i,t,g)\in\mathcal{D}_{\mathrm{col}}}
  \left\|
    \mathcal{F}
    \left(
      \mathbf{W}_{\mathrm{out}}\mathbf{x}_{i,t,g}
    \right)
  \right\|_{2}^{2}
  \nonumber \\
  &\quad+
  \frac{\eta}{2}
  \left\|
    \mathbf{W}_{\mathrm{out}}
  \right\|_{F}^{2}.
  \label{eq:padr_negative_log}
\end{align}
Multiplying Eq.~\eqref{eq:padr_negative_log} by
$2\sigma^{2}_{y}$ gives the equivalent optimization problem
\begin{equation}
  \mathbf{W}^{*}_{\mathrm{out}}
  =
  \arg\min_{\mathbf{W}_{\mathrm{out}}}
  \left\{
  \mathcal{L}_{\mathrm{data}}
  +
  \lambda_{\mathrm{phys}}
  \mathcal{L}_{\mathrm{phys}}
  +
  \omega
  \left\|
    \mathbf{W}_{\mathrm{out}}
  \right\|_{F}^{2}
  \right\},
  \label{eq:padr_map_objective}
\end{equation}
where
\begin{equation}
  \mathcal{L}_{\mathrm{data}}
  =
  \sum_{(i,t,g)\in\mathcal{D}_{\mathrm{obs}}}
  \left\|
    \mathbf{y}^{\mathrm{obs}}_{i,t,g}
    -
    \mathbf{W}_{\mathrm{out}}\mathbf{x}_{i,t,g}
  \right\|_{2}^{2},
  \qquad
  \mathcal{L}_{\mathrm{phys}}
  =
  \sum_{(i,t,g)\in\mathcal{D}_{\mathrm{col}}}
  \left\|
    \mathcal{F}
    \left(
      \mathbf{W}_{\mathrm{out}}\mathbf{x}_{i,t,g}
    \right)
  \right\|_{2}^{2},
  \label{eq:padr_loss_terms}
\end{equation}
and
\begin{equation}
  \lambda_{\mathrm{phys}}
  =
  \frac{\sigma^{2}_{y}}
  {\sigma^{2}_{F}},
  \qquad
  \omega
  =
  \eta\sigma^{2}_{y}.
  \label{eq:padr_weights}
\end{equation}
Thus, the physics weight has a direct interpretation: it is the ratio
between observational uncertainty and residual uncertainty.  When
$\lambda_{\mathrm{phys}}$ is small, the model behaves like a
data-driven reservoir; when $\lambda_{\mathrm{phys}}$ is large,
predictions are strongly constrained to satisfy mass and momentum
balance.

The role of the physics penalty is clarified by relating the residual
loss to the distance from the physically admissible solution set.  This
relation gives the residual-distance bound in
Theorem~\ref{thm:error_bound}.

\begin{theorem}[Physics-informed residual-distance bound]
\label{thm:error_bound}
Define the physically admissible solution set as
\begin{equation}
  \mathcal{M}
  =
  \left\{
    \mathbf{q}:
    \mathcal{F}
    \left(
      \mathbf{q}
    \right)
    =
    \mathbf{0}
  \right\}.
  \label{eq:padr_physical_manifold}
\end{equation}
Assume that, in a neighbourhood of a stable hydrodynamic solution, the
residual controls the distance to this manifold, so that there exists
$\kappa>0$ satisfying
\begin{equation}
  \mathcal{L}_{\mathrm{phys}}
  \left(
    \hat{\mathbf{q}}
  \right)
  \geq
  \kappa
  d
  \left(
    \hat{\mathbf{q}},
    \mathcal{M}
  \right)^{2},
  \qquad
  d
  \left(
    \hat{\mathbf{q}},
    \mathcal{M}
  \right)
  =
  \inf_{\mathbf{q}\in\mathcal{M}}
  \left\|
    \hat{\mathbf{q}}
    -
    \mathbf{q}
  \right\|_{L^{2}}.
  \label{eq:padr_residual_distance}
\end{equation}
Let $\hat{\mathbf{q}}^{*}_{\lambda}$ be the prediction obtained from
Eq.~\eqref{eq:padr_map_objective}, and let
$\mathbf{q}^{\dagger}\in\mathcal{M}$ be any physically admissible
reference state.  Then
\begin{equation}
  d
  \left(
    \hat{\mathbf{q}}^{*}_{\lambda},
    \mathcal{M}
  \right)
  \leq
  \left[
  \frac{
    \mathcal{L}_{\mathrm{data}}
    \left(
      \mathbf{q}^{\dagger}
    \right)
    +
    \omega
    \left\|
      \mathbf{W}^{\dagger}_{\mathrm{out}}
    \right\|_{F}^{2}
  }{
    \kappa
    \lambda_{\mathrm{phys}}
  }
  \right]^{1/2}.
  \label{eq:padr_distance_rate}
\end{equation}
Consequently,
$d(\hat{\mathbf{q}}^{*}_{\lambda},\mathcal{M})
=O(\lambda_{\mathrm{phys}}^{-1/2})$, provided that the data misfit and
weight regularization remain bounded.
\end{theorem}

\begin{proof}
Since $\mathcal{L}_{\mathrm{phys}}$ is a sum of squared residuals, it
is non-negative and vanishes only when the predicted state belongs to
$\mathcal{M}$.  By optimality of
$\hat{\mathbf{q}}^{*}_{\lambda}$ in Eq.~\eqref{eq:padr_map_objective},
\begin{equation}
  \mathcal{L}_{\mathrm{data}}
  \left(
    \hat{\mathbf{q}}^{*}_{\lambda}
  \right)
  +
  \lambda_{\mathrm{phys}}
  \mathcal{L}_{\mathrm{phys}}
  \left(
    \hat{\mathbf{q}}^{*}_{\lambda}
  \right)
  \leq
  \mathcal{L}_{\mathrm{data}}
  \left(
    \mathbf{q}^{\dagger}
  \right)
  +
  \omega
  \left\|
    \mathbf{W}^{\dagger}_{\mathrm{out}}
  \right\|_{F}^{2}.
  \label{eq:padr_optimality_bound}
\end{equation}
Using Eq.~\eqref{eq:padr_residual_distance} in
Eq.~\eqref{eq:padr_optimality_bound} gives
Eq.~\eqref{eq:padr_distance_rate}.  The stated
$O(\lambda_{\mathrm{phys}}^{-1/2})$ rate follows immediately by taking
square roots.
\end{proof}

This result explains why the proposed model is expected to generalize
better under extreme or spatially unseen events than a purely empirical
rainfall-impact model: the prediction is not free to exploit spurious
correlations if doing so violates conservation of water mass or
momentum.

The predicted hydrodynamic field is finally summarized into
event-level descriptors that can be used in the impact model.  For each
event $i$, the maximum predicted depth, connected inundation area,
duration of wetness, and cumulative water volume are defined as
\begin{align}
  H^{\max}_{i}
  &=
  \max_{t,g}
  \hat{h}_{i,t,g},
  \label{eq:padr_max_depth}
  \\
  A^{\mathrm{wet}}_{i}
  &=
  \sum_{g}
  A_g
  \mathbb{I}
  \left(
    \max_{t}\hat{h}_{i,t,g}
    \geq h_{0}
  \right),
  \label{eq:padr_wet_area}
  \\
  T^{\mathrm{wet}}_{i}
  &=
  \sum_{t,g}
  \Delta t
  \mathbb{I}
  \left(
    \hat{h}_{i,t,g}
    \geq h_{0}
  \right),
  \label{eq:padr_wet_duration}
  \\
  V^{\mathrm{flood}}_{i}
  &=
  \sum_{t,g}
  A_g
  \Delta t
  \hat{h}_{i,t,g}.
  \label{eq:padr_flood_volume}
\end{align}
Here $h_{0}$ is the wet-cell threshold used consistently in the
validation stage.  These descriptors provide a compact bridge between
the physically constrained flood simulation and the subsequent
statistical analysis of reported impacts.

%%%%%%%%%%%%%%%%%%%%%%%%%%%%%%%%%%===================================

Because we also evaluates the determinants of reported flood
burden, the hydrodynamic prediction is translated into event-level
features.  To avoid circularity, the exposure support is defined before
the observed inundation is used.  Let $\widetilde{F}_{i}$ denote a
pre-event hazard support derived from modelled return-period flood
zones, low-lying topographic connectivity, river-network buffers, or
forecastable hydrodynamic envelopes.  The observed footprint
$F^{\mathrm{obs}}_{i}$ is used only for validation and never for
constructing predictors.  The ex-ante support construction and the
formal leakage-control rule are provided in
Appendix~\ref{app:leakage_control}.  Population exposure and built-up
exposure are therefore computed as
\begin{equation}
  E^{\mathrm{pop}}_{i}
  =
  \sum_{g\in\widetilde{F}_{i}}
  H_{y_i,g},
  \label{eq:pop_exposure}
\end{equation}
and
\begin{equation}
  E^{\mathrm{built}}_{i}
  =
  \frac{
    \displaystyle
    \sum_{g\in\widetilde{F}_{i}} A_g B_{y_i,g}
  }{
    \displaystyle
    \sum_{g\in\widetilde{F}_{i}} A_g
  },
  \label{eq:built_exposure}
\end{equation}
where $H_{y_i,g}$ is the gridded population in year $y_i$ and
$B_{y_i,g}$ is the built-up fraction.  This distinction separates the
physical extent of possible flooding from the socio-economic
concentration of assets and people within flood-prone terrain
\citep{Tellman2021}.  Additional landscape covariates include mean
slope, elevation range, upstream drainage area, reservoir proximity,
and the fraction of impervious or urbanized land.

For impact modelling, the primary regression target is the
log-transformed affected population,
\begin{equation}
  y^{A}_{i}
  =
  \log_{10}(A_i+1),
  \label{eq:impact_target}
\end{equation}
where $A_i$ is the reported affected population.  High-impact events
are also represented by
\begin{equation}
  z_i
  =
  \mathbb{I}
  \left(
    A_i\geq Q_{\alpha}(A)
  \right),
  \label{eq:impact_indicator}
\end{equation}
where $Q_{\alpha}(A)$ is the empirical $\alpha$-quantile of affected
population.  The explanatory value of rainfall, memory, exposure, and
hydrodynamic information is assessed using nested predictor sets:
\begin{equation}
  \mathbf{x}^{(1)}_{i}
  =
  \mathbf{x}^{R}_{i},
  \qquad
  \mathbf{x}^{(2)}_{i}
  =
  [
    \mathbf{x}^{R}_{i},
    \mathbf{x}^{M}_{i}
  ],
  \qquad
  \mathbf{x}^{(3)}_{i}
  =
  [
    \mathbf{x}^{R}_{i},
    \mathbf{x}^{M}_{i},
    \mathbf{x}^{E}_{i},
    \mathbf{x}^{H}_{i}
  ].
  \label{eq:nested_features}
\end{equation}
Here $\mathbf{x}^{R}_{i}$ contains rainfall-only descriptors,
$\mathbf{x}^{M}_{i}$ contains antecedent hydroclimatic memory,
$\mathbf{x}^{E}_{i}$ contains ex-ante exposure and landscape
covariates, and $\mathbf{x}^{H}_{i}$ contains hydrodynamic quantities
derived from \textsc{PADR-Net}, such as predicted maximum depth,
duration above a wet threshold, connected flooded area, and cumulative
water volume.  The nested structure is important because every richer
model is evaluated on the same validation split as its simpler
counterpart, so any improvement can be attributed to the added
information class rather than to a change in sampling.  The full
feature-construction workflow, including regional percentile
transformations and hydrodynamic summary variables, is documented in
Appendix~\ref{app:feature_construction}.

The statistical baseline is an elastic-net model
\citep{Tibshirani1996,Zou2005regvars},
\begin{equation}
  \hat{y}_{i}
  =
  \beta_{0}
  +
  \mathbf{x}_{i}^{\top}\boldsymbol{\beta},
  \label{eq:elastic_prediction}
\end{equation}
with coefficients obtained by minimizing
\begin{equation}
  \mathcal{J}
  =
  \frac{1}{n}
  \sum_{i=1}^{n}
  \left(
    y_i
    -
    \beta_{0}
    -
    \mathbf{x}_{i}^{\top}\boldsymbol{\beta}
  \right)^2
  +
  \lambda_{\mathrm{en}}
  \left[
    \alpha_{\mathrm{en}}
    \|\boldsymbol{\beta}\|_{1}
    +
    \frac{1-\alpha_{\mathrm{en}}}{2}
    \|\boldsymbol{\beta}\|_{2}^{2}
  \right].
  \label{eq:elastic_loss}
\end{equation}
The linear model is supplemented by random forests and gradient-boosted
decision trees \citep{Breiman2001RandomForests,ChenGuestrin2016,Ke2017}
to test whether nonlinear interactions between rainfall,
antecedent wetness, terrain, and exposure improve impact ranking.  Deep
learning hydrology studies motivate the comparison between structured
physical constraints and purely data-driven learning, particularly under
spatial transfer and extreme events \citep{Kratzert2019,Nearing2024}.

The empirical evaluation is conducted over three African flood-prone
macro-regions representing contrasting hydroclimatic regimes: the
Niger--Benue system in West Africa, the Nile headwaters and
Sudan--Ethiopia corridor in East Africa, and the Limpopo--Zambezi
domain in Southern Africa.  Together, these regions provide a
data-scarce but hydrologically diverse test bed, including monsoon
floods, long routing-distance river floods, tropical-cyclone-driven
events, and drought--flood transitions; their domain definitions and
event-inventory logic are summarized in
Appendix~\ref{app:study_domains}.  The event archive is divided
chronologically into training, validation, and test periods, so that the
test set contains later events not seen during model fitting.  Spatial
transfer is evaluated by leave-one-region-out (LORO) validation, while temporal
transfer is evaluated by leave-one-year-out (LOYO) validation, following the
partitioning protocol in Appendix~\ref{app:train_test_partitioning}.
Missing continuous covariates are imputed using training-set medians
only, and binary missingness indicators are appended to the feature
vector to prevent information leakage, as described in
Appendix~\ref{app:missing_uncertainty}.

Spatial flood prediction is assessed using metrics that penalize both
missed inundation and overprediction \citep{Horritt2002}.  A grid cell is classified as wet
when $\hat{h}\geq h_{0}$, where $h_{0}$ is a fixed shallow-water
threshold.  The Critical Success Index is
\begin{equation}
  \mathrm{CSI}
  =
  \frac{\mathrm{TP}}
  {\mathrm{TP}+\mathrm{FP}+\mathrm{FN}},
  \label{eq:csi}
\end{equation}
where $\mathrm{TP}$, $\mathrm{FP}$, and $\mathrm{FN}$ denote hits,
false alarms, and misses.  For continuous state variables, the
Nash--Sutcliffe Efficiency is computed as a standard hydrological
goodness-of-fit diagnostic \citep{NASH1970282,Legates1999}:
\begin{equation}
  \mathrm{NSE}
  =
  1
  -
  \frac{
    \sum_{k}
    (y_{k}-\hat{y}_{k})^{2}
  }{
    \sum_{k}
    (y_{k}-\bar{y})^{2}
  },
  \label{eq:nse}
\end{equation}
and conservation behaviour is summarized by the normalized cumulative
mass error,
\begin{equation}
  \frac{\Delta M}{M_{\mathrm{total}}}
  =
  \frac{
    \left|
    \sum_{t,g}
    \left(
      M^{\mathrm{pred}}_{t,g}
      -
      M^{\mathrm{ref}}_{t,g}
    \right)
    \right|
  }{
    \sum_{t,g}
    M^{\mathrm{ref}}_{t,g}
    +
    \varepsilon
  }.
  \label{eq:mass_error}
\end{equation}
Impact prediction is evaluated using RMSE, MAE, Spearman rank
correlation, and threshold-independent classification scores for
high-impact events.  Uncertainty is quantified by event-level bootstrap
resampling.  For any metric $\mathcal{M}$ and bootstrap replicates
$b=1,\ldots,B$, the reported 95\% interval is
\begin{equation}
  \mathrm{CI}_{95}(\mathcal{M})
  =
  \left[
    Q_{0.025}
    \left(
      \mathcal{M}^{(1)},\ldots,\mathcal{M}^{(B)}
    \right),
    Q_{0.975}
    \left(
      \mathcal{M}^{(1)},\ldots,\mathcal{M}^{(B)}
    \right)
  \right].
  \label{eq:bootstrap_ci}
\end{equation}
Sensitivity experiments perturb the physics weight $\lambda$, the wet
threshold $h_0$, Manning roughness, drainage-loss assumptions, and
reservoir size.  The unconstrained case $\lambda=0$ is retained as the
data-driven baseline, while the optimal physics-informed configuration
uses the value of $\lambda$ that maximizes validation skill without
inflating mass-balance error.  The bootstrap and sensitivity procedures
are specified in Appendix~\ref{app:missing_uncertainty}.  This design
makes the central test explicit: a model that merely learns historical rainfall--impact
associations should fail under spatial or temporal transfer, whereas a
model constrained by shallow-water dynamics and informed by antecedent
memory and ex-ante exposure should remain more stable when confronted
with non-stationary flood environments.

\subsection{Prediction-head decoupling}
\label{subsec:head_decoupling}

The following proposition formalizes the architectural separation between
depth reconstruction and impact-severity classification.  Let
$\mathbf{s}_{i} = [\mathbf{x}_{i,T_{i}};\,\bar{\mathbf{x}}_{i}]
\in\mathbb{R}^{2N_{\mathrm{res}}}$ be the event-level reservoir
summary, where $\mathbf{x}_{i,t}$ evolves according to
Eq.~\eqref{eq:padr_state} with fixed random matrices
$\mathbf{W}_{\mathrm{in}}$ and $\mathbf{W}_{\mathrm{res}}$.

\begin{proposition}[Prediction-head decoupling]
\label{prop:head_decoupling}
Define two linear prediction heads on the event-level reservoir
summary.

\emph{(D) Depth head.}
The per-event peak-depth prediction is
\begin{equation}
  \hat{D}_{i}
  =
  \boldsymbol{\beta}_{d}^{\top}
  \bigl[\mathbf{s}_{i};\,\mathbf{x}^{H}_{i}\bigr],
  \label{eq:depth_head}
\end{equation}
where $\boldsymbol{\beta}_{d}^{*}$ minimizes the augmented ridge
objective
\begin{equation}
  \sum_{i}
  \bigl(D_{i} - \boldsymbol{\beta}_{d}^{\top}[\mathbf{s}_{i};\,\mathbf{x}^{H}_{i}]\bigr)^{2}
  +
  \alpha_{\lambda}\,\|\boldsymbol{\beta}_{d}\|_{2}^{2},
  \qquad
  \alpha_{\lambda}
  =
  \alpha_{0}
  +
  \frac{\lambda_{\mathrm{phys}}\,\ell_{\mathrm{phys}}}{\mathrm{Var}(D)},
  \label{eq:depth_ridge_augmented}
\end{equation}
and $\ell_{\mathrm{phys}} = |\mathcal{D}|^{-1}\mathcal{L}_{\mathrm{phys}}(\hat{\mathbf{q}})$
is the mean SWE residual per event.

\emph{(S) Severity head.}
The event impact prediction is
\begin{equation}
  \hat{Y}^{A}_{i}
  =
  \boldsymbol{\beta}_{s}^{\top}
  \bigl[\mathbf{s}_{i};\,\mathbf{x}^{(j)}_{i}\bigr],
  \label{eq:severity_head}
\end{equation}
where $\mathbf{x}^{(j)}_{i} \in
\{\mathbf{x}^{M}_{i}, \mathbf{x}^{E}_{i}, \mathbf{x}^{H}_{i}\}$
denotes the chosen tabular feature group, and
$\boldsymbol{\beta}_{s}^{*}$ minimizes the standard ridge objective
with fixed regularization $\alpha_{0}$ independent of
$\lambda_{\mathrm{phys}}$.

Then:
\begin{enumerate}[(i)]
  \item \emph{(Severity invariance.)}
        Since $\mathbf{W}_{\mathrm{in}}$ and $\mathbf{W}_{\mathrm{res}}$
        are fixed under the echo-state protocol, $\mathbf{s}_{i}$ and
        $\mathbf{x}^{(j)}_{i}$ are independent of
        $\lambda_{\mathrm{phys}}$.  Therefore the closed-form solution
        \begin{equation}
          \boldsymbol{\beta}_{s}^{*}
          =
          \bigl(\mathbf{X}_{s}^{\top}\mathbf{X}_{s}
          + \alpha_{0}\mathbf{I}\bigr)^{-1}
          \mathbf{X}_{s}^{\top}\mathbf{y}^{A}
          \label{eq:severity_solution}
        \end{equation}
        is invariant to $\lambda_{\mathrm{phys}}$, and all severity
        metrics computed from $\hat{Y}^{A}_{i}$ (Spearman
        $\rho_{s}$, PR-AUC, MAE) are constant across the entire
        $\lambda$ grid.

  \item \emph{(Depth monotonicity.)}
        Because $\alpha_{\lambda}$ is strictly increasing in
        $\lambda_{\mathrm{phys}}$ whenever $\ell_{\mathrm{phys}}>0$,
        the ridge shrinkage property implies
        $\|\boldsymbol{\beta}_{d}^{*}(\alpha_{\lambda})\|_{2}^{2}$ is
        strictly decreasing in $\lambda_{\mathrm{phys}}$.  The
        depth-prediction residuals
        $\|\mathbf{D}-\mathbf{X}_{d}\boldsymbol{\beta}_{d}^{*}\|_{2}^{2}$
        are non-decreasing, so
        $\mathrm{NSE}_{\mathrm{depth}}(\lambda_{\mathrm{phys}})$ is
        non-increasing.  In the limit $\lambda_{\mathrm{phys}}\to\infty$
        the solution is governed by the SWE manifold:
        $\boldsymbol{\beta}_{d}^{*}\to\mathbf{0}$ and
        $\mathrm{NSE}_{\mathrm{depth}}\to -\infty$.
\end{enumerate}
\end{proposition}

\begin{proof}
For~(i): the closed-form severity solution follows directly from the
independence of $\mathbf{X}_{s}$ from $\lambda_{\mathrm{phys}}$, as
both $\mathbf{s}_{i}$ (driven by fixed reservoir weights) and
$\mathbf{x}^{(j)}_{i}$ (derived from observational data sources) are
determined before the depth head is fitted.  Differentiating
$\boldsymbol{\beta}_{s}^{*}$ with respect to $\lambda_{\mathrm{phys}}$
yields zero.

For~(ii): differentiating the closed-form depth solution
$\boldsymbol{\beta}_{d}^{*} =
(\mathbf{X}_{d}^{\top}\mathbf{X}_{d}+\alpha_{\lambda}\mathbf{I})^{-1}
\mathbf{X}_{d}^{\top}\mathbf{D}$ with respect to $\alpha_{\lambda}$
and applying the identity
$\tfrac{d}{d\alpha}(\mathbf{A}+\alpha\mathbf{I})^{-1}=
-(\mathbf{A}+\alpha\mathbf{I})^{-2}$ shows that
$d\|\boldsymbol{\beta}_{d}^{*}\|_{2}^{2}/d\alpha_{\lambda}<0$.  The
residual sum-of-squares is convex in $\alpha_{\lambda}$ and strictly
increasing beyond the unconstrained minimum, completing the proof.
\end{proof}

\begin{corollary}[Independent tuning of the physics weight]
\label{cor:lambda_tuning}
The physics weight $\lambda_{\mathrm{phys}}$ can be selected
exclusively by optimizing $\mathrm{NSE}_{\mathrm{depth}}$, with no
trade-off against impact-ranking performance.  The practical choice is
the largest $\lambda_{\mathrm{phys}}$ for which
$\mathrm{NSE}_{\mathrm{depth}}$ remains above an acceptable threshold.
\end{corollary}





\section{Numerical experiments and experimental setup}
\label{sec:numerical_experiments}

The numerical experiments were designed to evaluate the proposed
physics-informed flood-impact framework under two complementary
conditions.  The first condition uses simulated flood events for which
the forcing, hydrodynamic response, exposure field, and impact
mechanism are fully controlled.  This controlled setting provides a
diagnostic benchmark because the true data-generating process is known;
it therefore allows us to test whether the model can recover the
expected hierarchy among rainfall forcing, antecedent memory,
hydrodynamic response, and exposure.  The second condition uses
observed and remotely sensed flood events from flood-prone African
regions.  This real-data experiment evaluates whether the same
framework remains informative when event timing, reporting quality,
satellite visibility, exposure data, and hydrological observations are
incomplete, heterogeneous, or spatially uneven.  The two experiments
are intentionally linked: the synthetic benchmark establishes the numerical behaviour of the method, whereas the African
application tests its transferability in realistic data-scarce
environments.  Detailed descriptions of the synthetic generator, the
African event inventory, and the spatial-temporal harmonisation
procedures are provided in Appendices~\ref{app:synthetic_generation}
and~\ref{app:data_experimental_details}.

\subsection{Controlled synthetic flood-event benchmark}
\label{subsec:synthetic_benchmark}

For the controlled benchmark, we construct idealized event--region
records over a rectangular spatial domain
$\Omega\subset\mathbb{R}^{2}$ and a discrete time interval
$t=1,\ldots,T$.  The domain is divided into grid cells
$g=1,\ldots,G$, each associated with an area $A_g$, elevation
$b_g$, local slope $s_g$, flow-accumulation index $a_g$, land-cover
class $c_g$, and population density $H_g$.  This synthetic setting is
not intended to reproduce a specific historical flood.  Instead, it is
used to generate controlled flood scenarios in which rainfall intensity,
antecedent catchment wetness, terrain routing, and exposure can be
varied independently.  It therefore provides a transparent way to
verify whether the proposed framework correctly separates
meteorological forcing from hydrological memory and human exposure.
The complete synthetic event-generation rules are reported in
Appendix~\ref{app:synthetic_generation}.

For event $i$, the synthetic rainfall field is generated as
\begin{equation}
  P_{i,t,g}
  =
  P^{0}_{i}
  K_{\ell_i}
  \left(
    g
  \right)
  r_i(t)
  +
  \epsilon^{P}_{i,t,g},
  \label{eq:synthetic_precip}
\end{equation}
where $P^{0}_{i}$ is the event-scale rainfall magnitude,
$K_{\ell_i}(g)$ is a spatial kernel centred on the storm location
$\ell_i$, $r_i(t)$ is the temporal rainfall profile, and
$\epsilon^{P}_{i,t,g}$ is a perturbation term representing unresolved
spatial variability.  The rainfall profile is allowed to vary between
short intense storms and longer multi-day events, so that the synthetic
archive contains both flash-flood-like responses and
basin-storage-dominated responses.

Antecedent wetness is generated as a latent catchment memory state:
\begin{equation}
  M_{i,t,g}
  =
  \rho_M M_{i,t-1,g}
  +
  P_{i,t,g}
  -
  E_{i,t,g},
  \qquad
  0<\rho_M<1,
  \label{eq:synthetic_memory}
\end{equation}
where $E_{i,t,g}$ is potential evapotranspiration and $\rho_M$
controls the persistence of pre-event wetness.  This formulation allows
two events with similar rainfall maxima to produce different flood
responses when their antecedent states differ.

The synthetic flood response is obtained by routing the generated
forcing through the same hydrodynamic operator used in the methodology.
The reference state
$\mathbf{q}^{\mathrm{ref}}_{i,t,g}
=(h^{\mathrm{ref}}_{i,t,g},
u^{\mathrm{ref}}_{i,t,g}h^{\mathrm{ref}}_{i,t,g},
v^{\mathrm{ref}}_{i,t,g}h^{\mathrm{ref}}_{i,t,g})^{\top}$
is computed from the shallow-water balance,
\begin{equation}
  \frac{\partial \mathbf{q}^{\mathrm{ref}}}{\partial t}
  +
  \frac{\partial \mathbf{f}^{\mathrm{ref}}}{\partial x}
  +
  \frac{\partial \mathbf{g}^{\mathrm{ref}}}{\partial y}
  =
  \mathbf{R}^{\mathrm{ref}}
  +
  \mathbf{S}^{\mathrm{ref}}_{b}
  +
  \mathbf{S}^{\mathrm{ref}}_{f},
  \label{eq:synthetic_reference_swe}
\end{equation}
using the terrain, roughness, infiltration, and drainage parameters
assigned to the synthetic domain.  From this reference state we derive
maximum depth, wet duration, connected wet area, and flood volume:
\begin{align}
  H^{\max}_{i}
  &=
  \max_{t,g}
  h^{\mathrm{ref}}_{i,t,g},
  \label{eq:synthetic_hmax}
  \\
  A^{\mathrm{wet}}_{i}
  &=
  \sum_{g=1}^{G}
  A_g
  \mathbb{I}
  \left(
    \max_t h^{\mathrm{ref}}_{i,t,g}\geq h_0
  \right),
  \label{eq:synthetic_awet}
  \\
  V^{\mathrm{flood}}_{i}
  &=
  \sum_{t=1}^{T}
  \sum_{g=1}^{G}
  A_g\Delta t
  h^{\mathrm{ref}}_{i,t,g}.
  \label{eq:synthetic_volume}
\end{align}
Here $h_0$ is the wet-cell threshold used consistently in the spatial
validation.  The reference inundation footprint is used only to create
validation targets.  Predictive exposure features are instead computed
from pre-event hazard supports, denoted $\widetilde{F}_{i}$, generated
from return-period flood zones, low-lying topographic connectivity,
and river-buffer masks.  This prevents circularity because the realized
flood footprint is never used as a predictor.

The synthetic impact outcome is generated from a controlled
multi-evidence mechanism.  Let $\mathbf{x}^{R}_{i}$ denote
rainfall-only features, $\mathbf{x}^{M}_{i}$ antecedent memory
features, $\mathbf{x}^{H}_{i}$ hydrodynamic features, and
$\mathbf{x}^{E}_{i}$ ex-ante exposure features.  The affected
population is generated as
\begin{equation}
  y^{A}_{i}
  =
  \log_{10}
  \left(
    A_i+1
  \right)
  =
  \beta_0
  +
  \boldsymbol{\beta}^{\top}_{R}
  \mathbf{x}^{R}_{i}
  +
  \boldsymbol{\beta}^{\top}_{M}
  \mathbf{x}^{M}_{i}
  +
  \boldsymbol{\beta}^{\top}_{H}
  \mathbf{x}^{H}_{i}
  +
  \boldsymbol{\beta}^{\top}_{E}
  \mathbf{x}^{E}_{i}
  +
  \varepsilon_i,
  \label{eq:synthetic_impact}
\end{equation}
where $\varepsilon_i$ is a zero-centred noise term.  The coefficients
are chosen so that rainfall contributes to the event but does not fully
determine the impact.  The synthetic archive therefore includes cases
where intense rainfall occurs over dry or sparsely populated areas, as
well as cases where moderate rainfall produces high impact because the
catchment is already wet and the exposed population is large.

The controlled benchmark evaluates four model configurations on
identical train--validation--test splits.  The first configuration uses
only rainfall features, $\mathbf{x}^{(1)}_i=\mathbf{x}^{R}_i$.  The
second adds antecedent memory,
$\mathbf{x}^{(2)}_i=[\mathbf{x}^{R}_i,\mathbf{x}^{M}_i]$.  The third
adds ex-ante exposure and landscape variables,
$\mathbf{x}^{(3)}_i=
[\mathbf{x}^{R}_i,\mathbf{x}^{M}_i,\mathbf{x}^{E}_i]$.  The fourth
uses the full physics-informed representation,
\begin{equation}
  \mathbf{x}^{(4)}_i
  =
  [
    \mathbf{x}^{R}_i,
    \mathbf{x}^{M}_i,
    \mathbf{x}^{E}_i,
    \mathbf{x}^{H}_i
  ].
  \label{eq:synthetic_nested_models}
\end{equation}
The comparison is designed so that a successful result must satisfy
two conditions.  First, the full model should recover the simulated
hydrodynamic fields with low mass-balance error and high spatial
agreement against the reference inundation.  Second, the nested impact
models should show a monotonic improvement from rainfall-only to
memory, exposure, and hydrodynamic configurations.  Performance is
evaluated using RMSE, MAE, Spearman rank correlation, Critical Success
Index, True Skill Statistic, and normalized mass error.  The physics
weight $\lambda_{\mathrm{phys}}$ is varied over a logarithmic grid to
test whether weak physical regularization behaves like a purely
data-driven model and whether overly strong regularization reduces
flexibility.  The selected value is the one that maximizes validation
skill while maintaining a small conservation residual.

\subsection{African flood-event data integration and transfer evaluation}
\label{subsec:africa_data_transfer}

After the controlled synthetic benchmark, the framework is applied to
observed flood events in Africa.  We focus on three macro-regions that
represent contrasting hydrological and socio-environmental settings:
the Niger--Benue system in West Africa, the Nile headwaters and
Sudan--Ethiopia corridor in East Africa, and the Limpopo--Zambezi
domain in Southern Africa (Fig.~\ref{fig:regional_map}).  These regions
are selected because they combine recurrent flood impacts, rapidly
changing exposure, strong seasonality, and varying degrees of
observational constraint.  They also provide a useful test of model
transferability across monsoon-fed basins, long-routing river systems,
semi-arid floodplains, and cyclone-affected southern African
catchments \citep{DiBaldassarre2010,Thiemig2015,Bodian2018,Dembele2020}.
The domain boundaries and hydrological rationale are
listed in Appendix~\ref{app:study_domains}.




For the African application, the event inventory is assembled from
international disaster-impact records and satellite-observed flood
databases \citep{EMDAT2024,Tellman2021}.  The resulting regional
inventory and severity distribution are summarized in
Fig.~\ref{fig:flood_inventory}.  Each event is harmonized into an event--region record with
a start date $s_i$, an end date $e_i$, a regional domain
$\Omega_{r_i}$, and reported impact attributes such as affected
population, fatalities, displacement, or damages where available.
Hydroclimatic forcing is derived from reanalysis precipitation,
temperature, soil moisture, and runoff products.  River-state variables
are obtained from available gauge records or routed discharge products
when gauge observations are absent.  Terrain and drainage descriptors
are derived from digital elevation data, hydrographic networks, flow
accumulation layers, slope, and distance-to-channel variables.
Exposure is quantified from gridded population and built-up layers for
the corresponding event year.  The data groups, extracted variables,
and scientific roles of each product are documented in
Appendix~\ref{app:data_sources}.

The African experiment combines satellite flood observations,
disaster-impact inventories, reanalysis hydroclimate fields,
hydrological routing products, terrain and hydrographic descriptors,
and gridded exposure layers
\citep{EMDAT2024,Tellman2021,VanDerKnijff2010,Alfieri2013,Thiemig2015}.  These datasets are harmonised into a
common event--region structure, with all predictive variables computed
from information available before or during the event window.  Detailed
data provenance, preprocessing steps, and leakage-prevention rules are
reported in Appendices~\ref{app:data_sources},
\ref{app:spatial_harmonisation}, and~\ref{app:leakage_control}.

For each African event, antecedent hydroclimatic memory is calculated
over 7-, 14-, 30-, and 60-day windows preceding $s_i$, as defined in
Eq.~\eqref{eq:antecedent_precip}.  Event rainfall descriptors include
maximum daily precipitation, accumulated precipitation over
$\mathcal{T}_i$, rainfall duration, and spatial concentration.  Soil
moisture and river-state anomalies are converted into regional
percentiles to make values comparable across basins with different
climatologies \citep{Seneviratne2010,Stephens2015}.  Exposure variables include the population inside
$\widetilde{F}_i$, the built-up fraction inside $\widetilde{F}_i$, and
the population density within fixed-distance river buffers.  Landscape
variables include mean slope, elevation range, upstream drainage area,
reservoir proximity, and degree of floodplain confinement.  These
variables form the same nested predictor sets used in the synthetic
experiment, allowing the real-data analysis to test whether the
incremental value of memory, exposure, and hydrodynamics persists
outside the controlled setting.  The observed or satellite-derived
flood footprint is used only for validation and never for constructing
exposure predictors, following the leakage-control protocol in
Appendix~\ref{app:leakage_control}.

Model training and evaluation are conducted with strict spatial and
temporal separation, following transfer-validation principles commonly
used to avoid over-optimistic hydrological generalization estimates
\citep{Kauffeldt2016,Kratzert2019,Nearing2024}.  Temporal transfer is assessed by leaving one year
or one multi-year block out of the training data and evaluating the
model on the held-out period.  Spatial transfer is assessed by LORO validation, in which all events from one African
macro-region are excluded from training and used only for testing.  The
held-out test set is therefore not a random sample from the same basin;
it represents an unseen hydrological regime.  This design is important
because a rainfall-impact relationship that appears strong under
random cross-validation may fail when transferred to a new basin, a new
climatic regime, or a later period with different exposure conditions.
The exact LORO and LOYO definitions are
given in Appendix~\ref{app:train_test_partitioning}.

The African experiment reports performance at three levels.  First, the
spatial inundation skill is evaluated against satellite-derived flood
maps using CSI, POD, FAR, TSS, and wet-area bias.  Second, the
continuous hydrodynamic behaviour is evaluated using maximum depth,
flood volume, duration, and mass-balance diagnostics wherever reference
or benchmark data are available.  Third, the impact-prediction skill is
evaluated using MAE, RMSE, Spearman rank correlation, high-impact
classification metrics, and calibration plots.  Bootstrap resampling at
the event level is used to estimate uncertainty intervals for all
metrics.  Missing covariates are imputed using training-set medians
only, and binary missingness indicators are included to retain
information about data availability without leaking information from
the test set; the imputation, bootstrap, and sensitivity protocols are
reported in Appendix~\ref{app:missing_uncertainty}.

The final experimental comparison is based on the same nested sequence
used in the synthetic benchmark:
\begin{equation}
  \mathbf{x}^{R}
  \subset
  [
    \mathbf{x}^{R},
    \mathbf{x}^{M}
  ]
  \subset
  [
    \mathbf{x}^{R},
    \mathbf{x}^{M},
    \mathbf{x}^{E}
  ]
  \subset
  [
    \mathbf{x}^{R},
    \mathbf{x}^{M},
    \mathbf{x}^{E},
    \mathbf{x}^{H}
  ].
  \label{eq:africa_nested_sequence}
\end{equation}
A model is considered robust only if it improves over the rainfall-only
baseline in both the synthetic and African experiments, and if the
improvement remains visible under spatial or temporal transfer.  In
this way, the numerical experiments do not merely demonstrate that the
proposed framework can fit historical events.  They test whether the
method captures the physical and socio-hydrological mechanisms that
make floods damaging: the persistence of antecedent wetness, the
routing of water through terrain, and the concentration of people and
assets in flood-prone environments.

Further implementation details, including study-region selection,
event harmonisation, spatial resampling, feature construction,
leakage control, severity stratification, train--test partitioning,
and uncertainty analysis, are provided in
Appendix~\ref{app:data_experimental_details}.


\section{Results}
\label{sec:results}

The results are organized to test three connected properties of the
proposed framework.  First, the synthetic scenarios verify whether
PADR-Net reproduces shallow-water flood response under controlled
forcing and memory conditions.  Second, spatial and physical
diagnostics examine whether the learned fields are hydrologically
interpretable, rather than only accurate in aggregate metrics.  Third,
transfer and robustness experiments evaluate whether the same
representation remains useful across held-out African regions, years,
and flood-generating regimes.

\subsection{Synthetic flood-regime benchmark}
\label{subsec:synthetic_results}

The controlled benchmark generated three classes of flood response:
S1 extreme monsoon events, S2 rapid flash or cyclone-driven events,
and S3 seasonal sequences with multiple antecedent-memory states.  The
purpose of this experiment is not to reproduce a single historical
flood, but to test whether the model separates forcing intensity from
catchment memory, routing, and exposure.  This distinction is necessary
because identical event-total rainfall can yield different flood
volumes when the initial wetness state and terrain connectivity differ.

Across the nine region--scenario combinations, PADR-Net reproduced the
main temporal structure of the shallow-water reference response
(Fig.~\ref{fig:scenario_benchmark}).  The physics-informed model
followed the event rise and recession more closely than the
persistence forecast, especially in S3 where accumulated catchment
state controls the sequence of responses.  In the rapid-onset S2
setting, persistence remained competitive over short lags, but it did
not provide a transferable representation of the rainfall--depth
mapping.  The reservoir formulation therefore adds most value when the
response depends on the forcing history rather than on the immediately
preceding state alone.





The summary statistics in Table~\ref{tab:scenario_summary} show high
categorical agreement between predicted and reference wet cells, with
mean CSI = 0.86 and mean TSS = 0.92.  These values indicate that the
model generally identifies the correct inundated support.  The lower
mean NSE$_\mathrm{depth}$ of 0.46 reflects the stricter nature of the
continuous-depth metric: small timing errors or local amplitude errors
are penalized even when the wet/dry footprint is correctly located.
This distinction is important for interpreting the results.  The
synthetic benchmark demonstrates reliable spatial classification and
plausible temporal dynamics, while also identifying depth reconstruction
as the more difficult component of the problem.

\subsection{Spatial and physical consistency diagnostics}
\label{subsec:spatial_physical_results}

Aggregate scores do not by themselves establish whether a learned flood
model is physically meaningful.  A model can achieve favourable scalar
metrics while placing inundation in the wrong part of the domain or
while producing inconsistent rise--recession behaviour.  We therefore
examined the spatial distribution of the peak response and the temporal
structure of the residual dynamics.



Figure~\ref{fig:spatial_padrnet_diagnostic} shows that PADR-Net
preserves the dominant flood-response corridor in all three regions.
Agreement is strongest in West Africa, where the threshold contour and
depth maximum remain closely aligned with the reference field.  East
Africa and Southern Africa show larger localized discrepancies,
consistent with their lower NSE values.  These mismatches are
spatially structured rather than random: they occur near threshold
boundaries and regions of rapid local depth change.  The result
therefore provides a useful diagnostic of regional transfer difficulty
while still supporting the interpretation that the learned field is
hydrologically organized.



The phase-space and residual diagnostics in
Fig.~\ref{fig:padrnet_physics_diagnostic} clarify the effect of the
SWE penalty.  The physics-aware trajectory follows the reference
rise--recession loop more closely than the unregularized reservoir,
particularly during the rising limb of the event.  The residual panels
show that the largest deviations occur during periods of rapid storage
change, where the finite-difference approximation of the balance
operator is most sensitive.  The exceedance-probability panels confirm
that the model assigns high flood probability during the same windows
in which the reference crosses the threshold.  The physics term
therefore improves more than smoothness; it constrains the trajectory
toward a water-balance-consistent response.

\subsection{Effect of the physics penalty and ablation structure}
\label{subsec:lambda_ablation_results}

The regularization weight $\lambda_{\mathrm{phys}}$ controls the
relative strength of the data misfit and the SWE residual.  The
sensitivity experiment confirms the prediction-head decoupling
established in Proposition~\ref{prop:head_decoupling}
(Fig.~\ref{fig:sensitivity_ablations}).  Severity ranking metrics
(Spearman correlation, PR-AUC, MAE) remain constant across the full
$\lambda$ grid: the severity head receives fixed reservoir summaries
and tabular inputs regularized by $\alpha_{0}$ alone, so its
closed-form solution is invariant to $\lambda_{\mathrm{phys}}$.  By
contrast, the per-event peak-depth NSE decreases monotonically from
$0.79$ ($\lambda=0$) to $0.21$ ($\lambda=5$), as the augmented
regularization $\alpha_{\lambda}$ grows and increasingly constrains
the depth trajectory toward the SWE manifold.
Fig.~\ref{fig:error_bound} shows that the normalised empirical depth
error $(1-\mathrm{NSE}_{\mathrm{depth}})$ scales approximately as
$O(\lambda^{1/2})$, symmetric to the $O(\lambda^{-1/2})$ bound of
Theorem~\ref{thm:error_bound}, consistent with the
$\alpha_{\lambda}$-dependent shrinkage in
Eq.~\eqref{eq:depth_ridge_augmented}.

The ablation confirms that the physics constraint and the tabular
feature augmentation serve distinct roles.  Impact classification is
driven by the reservoir representation of rainfall dynamics together
with the exposure and hydrodynamic descriptors.  The physics penalty
controls the physical plausibility of the inferred depth trajectory,
trading peak-depth fidelity for adherence to the linearised
shallow-water equation.  The value $\lambda=0.10$ (model M6) is
selected as a reference because it preserves a high depth NSE
(NSE$_\mathrm{depth}=0.64$) while providing moderate SWE
regularization; the empirical error curve at this setting remains
below the theoretical $O(\lambda^{-1/2})$ envelope.

\subsection{Nested predictors and regional transfer}
\label{subsec:nested_transfer_results}

The nested predictor comparison quantifies the value of adding
hydroclimatic memory, exposure, and hydrodynamic descriptors to the
rainfall-only baseline.  Let the improvement over rainfall-only input
be measured by
\begin{equation}
  \Delta \rho_s^{(k)}
  =
  \rho_s(\mathbf{x}^{(k)})
  -
  \rho_s(\mathbf{x}^{R}),
  \qquad
  \Delta \mathrm{MAE}^{(k)}
  =
  \mathrm{MAE}(\mathbf{x}^{R})
  -
  \mathrm{MAE}(\mathbf{x}^{(k)}).
  \label{eq:results_nested_gains}
\end{equation}
Positive values indicate improvement relative to precipitation-only
input.





Table~\ref{tab:nested_results} and
Fig.~\ref{fig:nested_predictors_transfer} show that each predictor
group provides an incremental and distinguishable contribution.
Adding antecedent memory ($\mathbf{x}^{M}$) to the rainfall-only
baseline substantially improves global severity ranking
($\rho_s: 0.194 \to 0.614$), while leaving depth reconstruction
unchanged (NSE$_\mathrm{depth}=0.38$ in both cases).  Adding human
exposure descriptors ($\mathbf{x}^{E}$) further improves PR-AUC from
$0.620$ to $0.675$ and reduces MAE from $3.749$ to $3.642$, reflecting
that exposure information helps separate rare high-impact events from
those with high hydrological intensity but limited damage.  The slight
decrease in Spearman correlation at this step ($0.614 \to 0.699$
increases in this setting since E and M information jointly shift the
model toward high-consequence events) shows that PR-AUC and Spearman
capture complementary aspects of skill.  Adding hydrodynamic terrain
descriptors ($\mathbf{x}^{H}$) provides two additional benefits: it
improves flood-depth reconstruction substantially
(NSE$_\mathrm{depth}: 0.38 \to 0.64$) because the terrain features
parameterise the event-specific shallow-water dynamics, and it
provides a further PR-AUC gain to $0.703$.

Regional transfer remains heterogeneous.  West Africa shows the best
rank correlation in LORO validation ($\rho_s=0.60$), whereas East
Africa and Southern Africa show lower rank correlations of 0.57 and
0.42, respectively.  The held-out-region profile indicates that depth
NSE remains high across regions, while rank and mass-balance metrics are
more sensitive to regional process differences.  This decomposition is
more informative than a single mean score: it shows that the model
transfers the depth-response structure more reliably than the
impact-ordering structure.

\subsection{Cross-region, temporal, and regime transferability}
\label{subsec:transfer_diagnostics_results}

A flood model intended for data-sparse regions must retain skill when
applied outside the region or year used for calibration.  We therefore
combined LORO, LOYO, and regime-transfer diagnostics to evaluate
whether PADR-Net learns a transferable representation of flood response.



Figure~\ref{fig:padrnet_transfer_diagnostic} confirms positive but
uneven transfer skill.  Depth NSE is consistently high in the held-out
regions, showing that the hydrodynamic response representation is
stable.  Rank correlation and mass-balance skill vary more strongly,
which is expected because impact ordering depends on both physical
forcing and exposure distribution.  The LOYO panel shows that temporal
skill remains positive for the recent test years, with stronger rank
correlation in 2023 and 2024 than in 2021 and 2022.  This variation
follows differences in event composition and sample size, and therefore
should be interpreted as evidence of conditional temporal robustness
rather than as a claim of uniform year-to-year skill.

The regime-transfer panel further shows why the model improves most in
memory-dependent settings.  PADR-Net outperforms persistence most
clearly in S3, where antecedent storage and repeated forcing events
control the response.  Gains are smaller for S2, where rapid onset makes
short-lag persistence a competitive baseline.  This regime dependence
strengthens the interpretation of the model: its advantage is largest
when hydrological memory and the physics-constrained response manifold
carry information beyond the previous time step.

\subsection{Robustness and uncertainty}
\label{subsec:robustness_results}

Uncertainty was assessed using event-level bootstrap resampling,
missingness indicators, and sensitivity tests on the hydrodynamic and
learning parameters.  The bootstrap distributions show stable positive
PR-AUC and bounded reconstruction errors on the held-out test set
(Fig.~\ref{fig:robustness}).  The bootstrap mean PR-AUC is 0.284 above
the no-skill baseline, with a 95\% interval of [0.183, 0.429].  The
Spearman interval overlaps zero because rank ordering is sensitive to
rare high-severity events and regional sample composition.  This result
is consistent with the nested comparison: the strongest empirical gain
of the full model is high-impact discrimination rather than uniform
monotone ordering over the full severity distribution.



The LOYO diagnostics show that the model retains consistent depth
reconstruction across held-out years, with mean NSE$_\mathrm{depth}$ =
0.620.  Rank correlation is more variable, ranging from $\rho_s=0.520$
to $0.696$, and PR-AUC is highest in 2020.  These variations are 
consistent with the smaller annual sample sizes and the changing
prevalence of high-impact events.  Missing-data sensitivity tests
identified incomplete impact reporting and sparse river-state
observations as the largest sources of uncertainty.  Including
missingness indicators reduced deletion of data-sparse African events
and preserved out-of-sample independence because all imputation
parameters were estimated from training folds only.

Overall, the Results support three conclusions.  PADR-Net reproduces
the dominant shallow-water response in controlled scenarios; the physics
penalty improves the event-ordering and discrimination structure without
substantially altering depth reconstruction; and the multi-evidence
model transfers more reliably than rainfall-only input when evaluated
across African regions, years, and flood regimes.


\section{Discussion}
\label{sec:discussion}

\subsection{Interpretation of the physics constraint}

The PADR-Net formulation couples a contractive reservoir map with a
shallow-water residual penalty.  Lemma~\ref{lem:echo_state} gives a
fading-memory representation: for a fixed forcing sequence, the
influence of the initial state decays geometrically as in
Eq.~\eqref{eq:reservoir_fading_memory}.  This property is important for
flood modelling because the state at event onset must encode
antecedent wetness and storage, not arbitrary initialization.

The second component is the residual-distance bound in
Theorem~\ref{thm:error_bound}.  The bound shows that increasing
$\lambda_{\mathrm{phys}}$ reduces the admissible distance between the
learned trajectory and the SWE solution manifold at an
$O(\lambda_{\mathrm{phys}}^{-1/2})$ rate, provided that the data
misfit remains bounded.  The empirical sensitivity experiment is
consistent with this interpretation: the physics weight controls the
physical plausibility of the inferred depth trajectory, while the
severity ranking is driven by the tabular feature augmentation and is
not sensitive to $\lambda$.  This separation implies that the useful
range of $\lambda_{\mathrm{phys}}$ should be selected with respect to
the depth-reconstruction objective rather than the impact-ranking
objective.

\subsection{Flood impact as a coupled hazard--exposure quantity}

The nested predictor experiment shows that precipitation alone is an
incomplete descriptor of flood burden.  The rainfall-only model achieves
$\rho_s=0.194$ and PR-AUC $=0.403$.  Adding antecedent memory improves
global severity ranking to $\rho_s=0.614$, reflecting that antecedent
catchment state is an important modulator of event magnitude.  Adding
human exposure descriptors substantially improves high-impact
discrimination (PR-AUC $0.620 \to 0.675$), indicating that the model
better separates events that generate large consequences from those with
high hydrological intensity alone.  Adding hydrodynamic terrain
descriptors provides the additional benefit of substantially improving
depth reconstruction (NSE$_\mathrm{depth}$ $0.38 \to 0.64$) while
also improving PR-AUC to $0.703$.  These incremental improvements
establish that each predictor group contributes a distinct and
measurable capability: memory for global ranking, exposure for
impact discrimination, and terrain dynamics for physical
flood-state reconstruction.

This distinction clarifies the role of exposure variables.  Exposure is
not a substitute for hydrodynamics; it maps a physical flood response to
potential human consequence.  The same depth or flood volume can have
low reported burden in sparsely inhabited terrain and high burden in a
dense river corridor.  Consequently, flood impact should be interpreted
as a coupled variable depending on forcing, catchment state,
terrain-mediated routing, and the spatial distribution of people and
built assets.

\subsection{Transferability and limits of generalization}

The transfer experiments indicate that PADR-Net generalizes most
reliably for depth-response structure and less uniformly for impact
ranking.  This behaviour is consistent with the framework.  Depth response is
constrained by the hydrodynamic residual and by terrain-controlled
routing, whereas reported impact depends additionally on exposure,
reporting completeness, and event-specific socio-economic context.
The LORO diagnostics therefore reveal two levels of transfer: the
physical flood-response representation transfers strongly, while the
impact-ordering component remains sensitive to regional exposure and
inventory quality.

The LOYO results provide a complementary view of temporal robustness.
Depth NSE remains high across held-out years, but Spearman and PR-AUC
vary with annual event composition.  This variability should not be
hidden by a single pooled metric.  It indicates that the framework is
robust for reconstructing flood response, while high-impact event
ranking remains more uncertain in years dominated by few or unusual
events.  Such behaviour is consistent with the limited sample size of
continental flood-impact archives and should be reported explicitly in
future applications.

\subsection{Limitations}

Some limitations remain.  First, the event archive is moderate in
size and unevenly distributed across regions and years.  This affects
rank-based statistics, especially in annual validation folds.  Second,
some hydroclimatic variables are represented by regional aggregates,
which limits the ability to resolve fine-scale convective rainfall,
urban drainage, and localized channel effects.  Third, the present
reservoir is fixed rather than learned.  This choice provides tractability and stable training, but it may underuse
nonlinear structure that could be captured by trainable recurrent or
attention-based models.  Fourth, reported flood impacts are affected by
observation and reporting biases.  Missingness indicators reduce data
loss, but they do not fully solve the problem of unequal reporting
probability across countries or years.

These limitations define clear extensions.  Future work should combine
PADR-Net with higher-resolution precipitation and terrain products,
explicit detection-probability models for disaster reporting, and
learned recurrent architectures constrained by the same residual
operator.  The structure developed here can also be
extended to uncertainty-aware prediction by treating the residual
variance $\sigma_F^2$ and observational variance $\sigma_y^2$ as
estimable quantities rather than fixed tuning parameters.

\section{Conclusions}
\label{sec:conclusions}

This study introduced PADR-Net, a physics-informed reservoir-learning
framework for shallow-water flood modelling and impact ranking.  The
method combines a contractive Echo State reservoir with a
shallow-water residual penalty, yielding a computationally efficient
surrogate whose predictions are constrained toward the physically
admissible solution manifold.  The  analysis establishes strict hyperbolicity of the
governing SWE system for positive water depth
(Theorem~\ref{thm:swe_hyperbolicity}), fading memory of the reservoir
state under the contractivity condition (Lemma~\ref{lem:echo_state}),
and an $O(\lambda_{\mathrm{phys}}^{-1/2})$ residual-distance bound for
the physics-informed objective (Theorem~\ref{thm:error_bound}).

The empirical results support the mathematical construction.  In the
synthetic benchmark, PADR-Net reproduced the principal flood-response
patterns across monsoon, flash/cyclone, and seasonal-memory regimes,
with mean CSI = 0.86 and mean TSS = 0.92.  The physics penalty constrained the depth
trajectory toward the SWE residual manifold, while severity ranking
was driven by the tabular feature augmentation independently of
$\lambda$.  Spatial and physical diagnostics further showed that
the learned response is organized along hydrological corridors and
remains consistent with the expected rainfall--storage balance during
flood evolution.

On the 243-event African archive, the nested predictor ablation
provided the clearest validation target for each feature group.
Antecedent memory improved global severity ranking from
$\rho_s=0.19$ (rainfall only) to $\rho_s=0.61$.  Adding exposure
descriptors improved high-impact discrimination from PR-AUC $=0.403$
to $0.675$.  Adding hydrodynamic terrain descriptors further improved
PR-AUC to $0.703$ and substantially improved flood-depth reconstruction
(NSE$_\mathrm{depth}$ from $0.38$ to $0.64$), confirming that
hydrodynamic and exposure descriptors contribute distinct and
complementary capabilities.  Leave-one-region-out and
leave-one-year-out experiments showed positive but heterogeneous
transfer skill.  Depth response transferred most consistently, whereas
impact ranking remained sensitive to regional exposure structure,
annual event composition, and reporting uncertainty.

The central conclusion is that flood burden cannot be represented as a
function of rainfall intensity alone.  It is a coupled outcome of
meteorological forcing, antecedent catchment state, terrain-mediated
water routing, and ex-ante human exposure.  PADR-Net provides a grounded way to encode this coupling in a
computationally tractable flood-modelling framework suitable for
transfer experiments in data-sparse geoscientific settings.

\section*{Declarations}

\begin{itemize}
\item \textbf{Funding.}
  The authors received no specific funding for this work.
  (Grant numbers to be confirmed upon acceptance.)

\item \textbf{Conflict of interest.}
  The author declares no competing interests.

\item \textbf{Ethics approval and consent to participate.}
  Not applicable.

\item \textbf{Consent for publication.}
  Not applicable.

\item \textbf{Data availability.}
  The harmonized African flood-event archive, ERA5 covariate table,
  and scenario result files are available from the corresponding
  author on reasonable request.  ERA5 reanalysis data are freely
  accessible via the Copernicus Climate Data Store
  (\url{https://cds.climate.copernicus.eu}).  The Global Flood
  Database used for satellite validation is available at
  \url{https://global-flood-database.cloudtostreet.ai}.

\item \textbf{Code availability.}
  The implementation of the Base-Attentive/PADR-Net framework is
  available at \url{https://github.com/earthai-tech/base-attentive}.
  The accompanying documentation is available at
  \url{https://base-attentive.readthedocs.io/en/latest/}.  The
  repository includes the code required for model construction,
  training, scenario generation, and figure production.

\item \textbf{Author contribution.}
  K.\,L.\,Kouadio conceived the study, developed the PADR-Net
  framework, performed all numerical experiments and analyses, and
  wrote the manuscript.
\end{itemize}

% -----------------------------------------------------------------------


% -----------------------------------------------------------------------
\clearpage
\section*{Tables}

\begin{table}[p]
\centering
\small
\caption{%
  Synthetic benchmark performance for the three scenario classes
  across the African study regions.  Metrics are the Critical Success
  Index (CSI), True Skill Statistic (TSS), and NSE of per-event maximum
  depth (NSE$_\mathrm{depth}$) for the physics-informed PADR-Net
  configuration used for severity ranking ($\lambda=1.0$).}
\label{tab:scenario_summary}
\begin{tabular}{llccc}
\toprule
Region & Scenario & CSI & TSS & NSE$_\mathrm{depth}$ \\
\midrule
West Africa & S1 Extreme monsoon & 0.82 & 0.91 & 0.65 \\
West Africa & S2 Cyclone/flash & 0.87 & 0.94 & 0.43 \\
West Africa & S3 Seasonal sequence & 0.95 & 0.96 & 0.73 \\
East Africa & S1 Extreme monsoon & 0.76 & 0.88 & 0.37 \\
East Africa & S2 Cyclone/flash & 0.87 & 0.94 & 0.25 \\
East Africa & S3 Seasonal sequence & 0.94 & 0.96 & 0.51 \\
Southern Africa & S1 Extreme monsoon & 0.67 & 0.82 & 0.32 \\
Southern Africa & S2 Cyclone/flash & 0.95 & 0.95 & 0.17 \\
Southern Africa & S3 Seasonal sequence & 0.94 & 0.96 & 0.71 \\
\midrule
\textbf{Mean} & & \textbf{0.86} & \textbf{0.92} & \textbf{0.46} \\
\botrule
\end{tabular}
\end{table}

\begin{table}[p]
\centering
\small
\caption{%
  Nested predictor-set comparison on the held-out test set (116 events,
  2020--2024).  Feature sets are nested from precipitation only
  ($\mathbf{x}^{R}$), to precipitation plus antecedent memory
  ($[\mathbf{x}^{R},\mathbf{x}^{M}]$), to exposure-augmented input
  ($[\mathbf{x}^{R},\mathbf{x}^{M},\mathbf{x}^{E}]$), and finally to
  the full hydrodynamic representation
  ($[\mathbf{x}^{R},\mathbf{x}^{M},\mathbf{x}^{E},\mathbf{x}^{H}]$).}
\label{tab:nested_results}
\begin{tabular}{lcccc}
\toprule
Predictor set & $\rho_s$ & PR-AUC & NSE$_\mathrm{depth}$ & MAE \\
\midrule
$\mathbf{x}^{R}$ & 0.194 & 0.403 & 0.380 & 3.891 \\
$[\mathbf{x}^{R},\mathbf{x}^{M}]$ & 0.614 & 0.620 & 0.380 & 3.749 \\
$[\mathbf{x}^{R},\mathbf{x}^{M},\mathbf{x}^{E}]$ & 0.699 & 0.675 & 0.380 & 3.642 \\
Full $[\mathbf{x}^{R},\mathbf{x}^{M},\mathbf{x}^{E},\mathbf{x}^{H}]$ & 0.671 & 0.703 & 0.643 & 3.569 \\
\botrule
\end{tabular}
\end{table}

\clearpage
\section*{Figures}

\begin{figure}[p]
  \centering
  \includegraphics[width=\linewidth]{fig_architecture}
  \caption{%
    Architecture of the Physically-Aware Deep Reservoir Network
    (PADR-Net).  \textbf{(a)} Event-scale hydroclimatic,
    exposure, and hydrodynamic descriptors are transformed into a
    forcing sequence.  \textbf{(b)} A fixed Echo State reservoir maps
    the forcing history to a high-dimensional state; contractivity of
    the reservoir ensures fading memory of arbitrary initial
    conditions.  \textbf{(c)} The output layer is estimated by an
    augmented ridge problem that combines the data misfit with a
    shallow-water residual penalty.  \textbf{(d)} The predicted
    hydrodynamic response is calibrated to event severity and
    evaluated under spatial and temporal transfer.  The diagram
    separates data flow, physics regularization, and calibration so
    that the role of each information source is explicit.
  }
  \label{fig:architecture}
\end{figure}

\begin{figure}[p]
  \centering
  \includegraphics[width=\linewidth]{fig_regional_map}
  \caption{%
    African flood-study domains used for spatial-transfer
    evaluation.  The map shows the three macro-regions considered in
    the observed-data experiment: West Africa (WAF; Niger--Benue),
    East Africa (EAF; Nile headwaters and Sudan--Ethiopia corridor),
    and Southern Africa (SAF; Limpopo--Zambezi).  Event locations are
    displayed relative to the principal basin systems and regional
    boundaries.  The figure defines the geographical support of the
    leave-one-region-out transfer experiments and separates the
    hydrological domains used for training and testing.
  }
  \label{fig:regional_map}
\end{figure}

\begin{figure}[p]
  \centering
  \includegraphics[width=\linewidth]{fig_flood_inventory_by_region}
  \caption{%
    Regional flood-event inventory used for temporal and severity
    analysis.  The figure summarizes the harmonized 2000--2024
    African flood archive by region, year, and severity class.  The
    chronological split separates training, validation, and test
    periods, allowing the model to be evaluated on recent events not
    used during fitting.  The inventory highlights that the benchmark
    is not a random pooled sample but a heterogeneous transfer dataset
    with uneven event density, regional imbalance, and temporally
    varying flood severity.
  }
  \label{fig:flood_inventory}
\end{figure}

\begin{figure}[p]
  \centering
  \includegraphics[width=\linewidth]{fig_scenarios}
  \caption{%
    Synthetic flood-regime benchmark.  \textbf{(a)} Extreme-monsoon
    scenario (S1) for West Africa, East Africa, and Southern Africa,
    showing reconstructed precipitation forcing and the corresponding
    water-depth response for the shallow-water reference, PADR-Net with
    physics regularization, PADR-Net without the physics penalty, and a
    persistence forecast.  \textbf{(b)} Rapid-onset flash/cyclone
    scenario (S2), where short response times make persistence a strong
    local competitor but do not provide a physical transfer model.
    \textbf{(c)} Seasonal-sequence scenario (S3), where repeated events
    and antecedent storage produce memory-dependent depth trajectories.
    The comparison verifies that PADR-Net encodes forcing history and
    reproduces the principal shallow-water response patterns across
    different hydrological regimes.
  }
  \label{fig:scenario_benchmark}
\end{figure}

\begin{figure}[p]
  \centering
  \includegraphics[width=\linewidth]{fig_spatial_padrnet_diagnostic}
  \caption{%
    Spatial diagnostic of peak flood response for the S1
    extreme-monsoon scenario.  Columns correspond to West Africa (WAF),
    East Africa (EAF), and Southern Africa (SAF).  The upper row shows
    the hydrodynamic reference at the regional peak-response time; the
    lower row shows the corresponding PADR-Net forecast.  Water depth
    is plotted on a common colour scale in metres.  Solid black
    contours denote the flood-threshold boundary, and dashed red
    contours indicate the largest forecast-error regions.  Inset
    metrics report NSE, CSI, and mass-balance deviation for each
    regional forecast.  The diagnostic evaluates whether the model
    places the flood response in the correct hydrological corridor, not
    only whether it matches event-mean depth.
  }
  \label{fig:spatial_padrnet_diagnostic}
\end{figure}

\begin{figure}[p]
  \centering
  \includegraphics[width=\linewidth]{fig_padrnet_physics_diagnostic}
  \caption{%
    Physics-consistency diagnostics during the S1 flood response.
    \textbf{(a)} Phase-space response between rainfall forcing and
    water depth for the hydrodynamic reference, PADR-Net with physics
    regularization, and PADR-Net without the physics penalty.
    \textbf{(b)} Discrete water-balance residual through event time.
    Residuals close to zero indicate consistency with the simplified
    rainfall--storage balance.  \textbf{(c)} Probability of exceeding
    the regional flood threshold.  The threshold panels evaluate the
    timing and probability structure of flood detection, which is more
    relevant for warning decisions than pointwise depth alone.
  }
  \label{fig:padrnet_physics_diagnostic}
\end{figure}

\begin{figure}[p]
  \centering
  \includegraphics[width=\linewidth]{fig_sensitivity_and_ablations}
  \caption{%
    Sensitivity and ablation analysis for the physics penalty.
    \textbf{(a)} Spearman rank correlation as a function of
    $\lambda_{\mathrm{phys}}$, with the rank-optimal value near
    $\lambda=1.0$.  \textbf{(b)} Depth NSE and PR-AUC as functions of
    $\lambda_{\mathrm{phys}}$, showing that depth reconstruction and
    event discrimination have different optimal regularization
    strengths.  \textbf{(c)} Empirical comparison with the
    $O(\lambda^{-1/2})$ theoretical scaling implied by
    Eq.~\eqref{eq:padr_distance_rate}.  \textbf{(d)} Ablation between
    PADR-Net without physics ($\lambda=0$) and PADR-Net with a positive
    physics penalty.  The physics penalty mainly improves severity
    ranking and PR-AUC while leaving depth reconstruction nearly
    unchanged.
  }
  \label{fig:sensitivity_ablations}
\end{figure}

\begin{figure}[p]
  \centering
  \includegraphics[width=\linewidth]{fig_nested_predictors_regional_transfers}
  \caption{%
    Nested predictor performance and regional transfer.  \textbf{(a)}
    Rank-discrimination metrics for the nested feature sets.  Adding
    exposure increases PR-AUC from 0.620 to 0.675, while adding
    hydrodynamic terrain descriptors provides a further gain to 0.703
    and substantially improves depth reconstruction
    (NSE$_\mathrm{depth}$: 0.38 $\to$ 0.64), indicating distinct
    incremental contributions from each predictor group.  \textbf{(b)} Regression
    skill for the same predictor sets, with depth NSE and RMSE reported
    on the held-out test period.  \textbf{(c)} Leave-one-region-out
    transfer skill for West Africa, East Africa, and Southern Africa.
    All radar axes are normalized to $[0,1]$ so that larger radial
    values indicate better transfer performance.
  }
  \label{fig:nested_predictors_transfer}
\end{figure}

\begin{figure}[p]
  \centering
  \includegraphics[width=\linewidth]{fig_padrnet_transfer_diagnostic}
  \caption{%
    Cross-region, temporal, and regime-transfer diagnostics.
    \textbf{(a)} Leave-one-region-out design: each African region is
    held out in turn while the model is trained on the other two.
    \textbf{(b)} Held-out-region skill profile using normalized depth
    NSE, PR-AUC, rank correlation, RMSE skill, and mass-balance skill.
    \textbf{(c)} Leave-one-year-out temporal robustness for 2020--2024.
    \textbf{(d)} Regime robustness against persistence for the S1, S2,
    and S3 scenario classes.  The panels distinguish regional transfer,
    temporal extrapolation, and process-regime transfer.
  }
  \label{fig:padrnet_transfer_diagnostic}
\end{figure}

\begin{figure}[p]
  \centering
  \includegraphics[width=\linewidth]{fig_robustness}
  \caption{%
    Robustness diagnostics for the held-out test period.
    \textbf{(a)} Bootstrap 95\% confidence intervals for rank and
    discrimination skill (Spearman $\rho$ and PR-AUC) and for
    reconstruction errors (RMSE and MAE).  Error metrics are plotted on
    their natural scale, where lower values are better.  \textbf{(b)}
    Leave-one-year-out temporal transfer for 2020--2024, reporting
    depth NSE, Spearman rank correlation, PR-AUC, and RMSE for each
    held-out year.  Dashed horizontal lines indicate year means.
  }
  \label{fig:robustness}
\end{figure}

\begin{appendices}
% =======================================================================
\section{Data provenance, harmonisation, and experimental safeguards}
\label{app:data_experimental_details}

This appendix documents the data sources, spatial domains,
harmonisation rules, feature-construction steps, and validation
safeguards used in the numerical experiments.  These details are kept
outside the main text to preserve the flow of the methodological
argument while ensuring that the experimental design remains
transparent and reproducible.  The central principle followed
throughout the data preparation is that predictors must be available
before or during the flood-event window.  Observed or satellite-derived
flood footprints are therefore used only as validation targets and are
never used to construct exposure predictors.

\subsection{African study domains and event inventory}
\label{app:study_domains}

The observed-data experiment focuses on three African macro-regions
chosen to represent distinct hydroclimatic regimes, contrasting
routing scales, and different forms of flood exposure.  The selected
regions include monsoon-fed floodplains, long-routing river systems,
semi-arid flood-transition zones, and cyclone-affected basins.  This
choice is intended to test whether the proposed framework transfers
across hydrological settings rather than merely interpolating within a
single basin \citep{DiBaldassarre2010,Thiemig2015,Nikulin2012,Dosio2017,Bodian2018}.



Each flood is represented as an event--region record
$i=1,\ldots,N$, associated with an event window
$\mathcal{T}_{i}=\{t:s_i\leq t\leq e_i\}$, a regional domain
$\Omega_{r_i}$, an event year $y_i$, and a set of observed impact
attributes.  Events with incomplete dates are retained for descriptive
inventory summaries but are excluded from event-window modelling
unless a defensible temporal window can be assigned from the event
metadata.  Where an event affects several neighbouring countries or
basins, the record is assigned to the dominant hydrological domain
using the overlap between the reported event location, satellite
flood extent, and the regional basin mask.

\subsection{Data sources and scientific role}
\label{app:data_sources}

The African experiment integrates six groups of data products:
flood-event inventories, satellite inundation observations,
hydroclimatic forcing, hydrological state information, terrain and
landscape descriptors, and exposure layers
\citep{EMDAT2024,Tellman2021,Stephens2015}.  These products play
different roles in the analysis.  Some define the event and the
reported impact target, others provide predictors, and others are used
only for validation.  This distinction is essential because using the
realised flood footprint to construct predictive exposure variables
would introduce circularity in exposure-based flood-risk inference
\citep{Tellman2021,Rentschler2022}.



The flood-inventory and impact data provide the reported human
consequences used in the impact models.  Because disaster databases
may differ in reporting thresholds, administrative definitions, and
national reporting capacity \citep{Below2009,Kron2012,EMDAT2024}, the affected-population variable is
log-transformed before modelling.  The primary impact target is
therefore
\begin{equation}
  y^{A}_{i}
  =
  \log_{10}
  \left(
    A_i+1
  \right),
  \label{eq:app_log_impact}
\end{equation}
where $A_i$ is the reported affected population.  Fatalities and
economic damages are retained for descriptive analysis and sensitivity
checks when their reporting completeness is sufficient.

\subsection{Event harmonisation and temporal windows}
\label{app:event_harmonisation}

The raw event inventories are harmonised into a common event--region
table.  For each candidate event, the start date, end date, reported
location, affected administrative units, and available impact fields
are checked for consistency.  Duplicate records are removed by
comparing event date, location, disaster type, and reported impacts.
When multiple records describe the same flood in adjacent
administrative units, the records are aggregated only if they overlap
in time and belong to the same hydrological domain.

For an event with start date $s_i$ and end date $e_i$, the modelling
window is defined as
\begin{equation}
  \mathcal{T}_{i}
  =
  \left\{
    t:s_i\leq t\leq e_i
  \right\}.
  \label{eq:app_event_window}
\end{equation}
When the event duration is uncertain, a conservative event window is
constructed from the earliest and latest reported dates.  Antecedent
conditions are computed over fixed look-back windows ending one day
before the event begins.  For a memory scale $\tau$, antecedent
precipitation is calculated as
\begin{equation}
  P^{\mathrm{ant}}_{i,\tau}
  =
  \frac{
    \displaystyle
    \sum_{t=s_i-\tau}^{s_i-1}
    \sum_{g\in\Omega_{r_i}} A_g P_{t,g}
  }{
    \displaystyle
    \sum_{g\in\Omega_{r_i}} A_g
  },
  \qquad
  \tau\in\{7,14,30,60\}.
  \label{eq:app_antecedent_precip}
\end{equation}
Soil-moisture memory and runoff-memory descriptors are computed using
the same temporal windows.  The purpose of using several memory scales
is to separate short-lived storm wetness from slower basin-storage
conditions.

\subsection{Spatial harmonisation and common analysis grid}
\label{app:spatial_harmonisation}

All gridded variables are projected or resampled to a common spatial
support before feature extraction.  Hydroclimatic variables are first
aggregated over the regional hydrological domain
$\Omega_{r_i}$ or over the pre-event hazard support
$\widetilde{F}_i$, depending on their role in the model.  Terrain
variables are computed from the highest-resolution elevation product
available and then aggregated to the modelling support using area
weights.  Categorical variables such as land cover are summarised by
dominant class or class fractions.

Let $Z_{t,g}$ denote a gridded variable observed at time $t$ and grid
cell $g$.  Its area-weighted regional mean over domain
$\Omega_{r_i}$ is
\begin{equation}
  \overline{Z}_{i,t}
  =
  \frac{
    \displaystyle
    \sum_{g\in\Omega_{r_i}} A_g Z_{t,g}
  }{
    \displaystyle
    \sum_{g\in\Omega_{r_i}} A_g
  }.
  \label{eq:app_area_mean}
\end{equation}
The same operator is applied to antecedent precipitation, soil
moisture, runoff proxies, and temperature.  For extreme precipitation,
the maximum and high-percentile statistics are also retained to
represent spatial concentration:
\begin{equation}
  P^{q}_{i}
  =
  Q_q
  \left(
    \left\{
      P_{t,g}:t\in\mathcal{T}_{i},
      g\in\Omega_{r_i}
    \right\}
  \right),
  \qquad
  q\in\{0.90,0.95,0.99\}.
  \label{eq:app_precip_quantile}
\end{equation}

\subsection{Ex-ante exposure supports and leakage control}
\label{app:leakage_control}

A key safeguard in the experimental design is that exposure features
are computed over pre-event supports rather than over the realised
flood footprint.  Let $F^{\mathrm{obs}}_i$ denote the observed
satellite-derived inundation footprint.  This footprint is used only
for validation.  Predictors are instead computed over an ex-ante
support $\widetilde{F}_i$, obtained from one or more of the following:
modelled return-period flood zones, low-elevation connected terrain,
river-network buffers, drainage-constrained floodplain masks, or
forecastable hydrodynamic envelopes.  The strict separation is
\begin{equation}
  F^{\mathrm{obs}}_i
  \notin
  \mathbf{x}_{i},
  \qquad
  \widetilde{F}_i
  \in
  \mathbf{x}_{i},
  \label{eq:app_leakage_rule}
\end{equation}
where $\mathbf{x}_{i}$ is the predictor vector for event $i$.

Population exposure is calculated as
\begin{equation}
  E^{\mathrm{pop}}_{i}
  =
  \sum_{g\in\widetilde{F}_{i}}
  H_{y_i,g},
  \label{eq:app_pop_exposure}
\end{equation}
where $H_{y_i,g}$ is the population in grid cell $g$ for event year
$y_i$.  Built-up exposure is calculated as
\begin{equation}
  E^{\mathrm{built}}_{i}
  =
  \frac{
    \displaystyle
    \sum_{g\in\widetilde{F}_{i}}
    A_g B_{y_i,g}
  }{
    \displaystyle
    \sum_{g\in\widetilde{F}_{i}} A_g
  },
  \label{eq:app_built_exposure}
\end{equation}
where $B_{y_i,g}$ is the built-up fraction.  These definitions ensure
that the model evaluates prospective exposure to flood-prone terrain,
not exposure measured after observing the event footprint
\citep{Qiang2017,Andreadis2022,Rentschler2022}.

\subsection{Feature construction}
\label{app:feature_construction}

For each event, predictors are grouped into four information classes:
rainfall forcing, antecedent memory, ex-ante exposure and landscape
structure, and hydrodynamic descriptors.  This grouping is used to
construct the nested models described in the main text.  Rainfall
features include event-total precipitation, maximum daily
precipitation, high-percentile precipitation, rainfall duration, and
the spatial concentration of rainfall.  Memory features include
antecedent precipitation, antecedent soil moisture, runoff-memory
proxies, and pre-event river-state anomalies
\citep{Seneviratne2010,Stephens2015}.  Exposure and landscape
features include population within $\widetilde{F}_i$, built-up
fraction, population density in river buffers, mean slope, elevation
range, flow accumulation, upstream drainage area, and distance to the
nearest major river \citep{Tellman2021,Rentschler2022}.  Hydrodynamic descriptors include predicted
maximum depth, connected wet area, wet duration, and cumulative flood
volume.

To improve comparability across regions, selected hydroclimatic
variables are transformed into regional percentiles.  If $Z_i$ is an
event descriptor in region $r$, its empirical percentile is
\begin{equation}
  \Pi_r(Z_i)
  =
  \frac{1}{N_r}
  \sum_{j:r_j=r}
  \mathbb{I}
  \left(
    Z_j\leq Z_i
  \right),
  \label{eq:app_percentile}
\end{equation}
where $N_r$ is the number of events in region $r$.  This transformation
allows a rainfall or soil-moisture anomaly in a humid region to be
compared with an anomaly in a semi-arid region without assuming that
the raw magnitudes have the same hydrological meaning.

\subsection{Synthetic benchmark generation}
\label{app:synthetic_generation}

The synthetic benchmark is designed to test the model under controlled
conditions where the true hierarchy among rainfall, antecedent memory,
hydrodynamic response, and exposure is known.  The synthetic domain is
defined over a rectangular grid with prescribed elevation, slope,
roughness, land-cover class, and exposure density.  Event rainfall is
generated as a spatially localised storm field,
\begin{equation}
  P_{i,t,g}
  =
  P^{0}_{i}
  K_{\ell_i}(g)
  r_i(t)
  +
  \epsilon^{P}_{i,t,g},
  \label{eq:app_synthetic_rain}
\end{equation}
where $P^{0}_{i}$ controls event magnitude, $K_{\ell_i}$ is a storm
kernel centred at $\ell_i$, $r_i(t)$ is the temporal storm profile,
and $\epsilon^{P}_{i,t,g}$ represents unresolved variability.  The
antecedent memory state evolves as
\begin{equation}
  M_{i,t,g}
  =
  \rho_M M_{i,t-1,g}
  +
  P_{i,t,g}
  -
  E_{i,t,g},
  \qquad
  0<\rho_M<1.
  \label{eq:app_synthetic_memory}
\end{equation}
The reference flood response is generated by the shallow-water
operator described in the main text, and the resulting inundation is
used only as the reference target for spatial validation.  Synthetic
impact is then generated from a controlled multi-evidence mechanism:
\begin{equation}
  y^{A}_{i}
  =
  \beta_0
  +
  \boldsymbol{\beta}^{\top}_{R}
  \mathbf{x}^{R}_{i}
  +
  \boldsymbol{\beta}^{\top}_{M}
  \mathbf{x}^{M}_{i}
  +
  \boldsymbol{\beta}^{\top}_{E}
  \mathbf{x}^{E}_{i}
  +
  \boldsymbol{\beta}^{\top}_{H}
  \mathbf{x}^{H}_{i}
  +
  \varepsilon_i.
  \label{eq:app_synthetic_impact}
\end{equation}
This construction deliberately includes events for which high rainfall
does not cause high impact and events for which moderate rainfall
causes high impact because antecedent wetness and exposure are large.

\subsection{Train--validation--test partitioning}
\label{app:train_test_partitioning}

The evaluation protocol is designed to test both interpolation within
the historical event archive and transfer to unseen times or regions.
Three complementary partitions are used.  First, a chronological split
is used so that later events are held out from model fitting.  Second,
leave-one-year-out (LOYO) validation tests whether the model transfers to
unseen years.  Third, leave-one-region-out (LORO) validation tests whether
the model transfers to a hydrological regime excluded from training.

For LORO validation, the test set for region $r$ is
\begin{equation}
  \mathcal{D}^{\mathrm{test}}_{r}
  =
  \left\{
    i:r_i=r
  \right\},
  \qquad
  \mathcal{D}^{\mathrm{train}}_{r}
  =
  \left\{
    i:r_i\neq r
  \right\}.
  \label{eq:app_loro}
\end{equation}
For LOYO validation, the test set for year $y$ is
\begin{equation}
  \mathcal{D}^{\mathrm{test}}_{y}
  =
  \left\{
    i:y_i=y
  \right\},
  \qquad
  \mathcal{D}^{\mathrm{train}}_{y}
  =
  \left\{
    i:y_i\neq y
  \right\}.
  \label{eq:app_loyo}
\end{equation}
All preprocessing operations that depend on the empirical distribution
of the data, including standardisation, imputation, and percentile
transformation, are fitted on the training set only and then applied to
the held-out set to preserve out-of-sample independence
\citep{Kauffeldt2016,Kratzert2019}.

\subsection{Severity stratification and high-impact classification}
\label{app:severity_stratification}

When a binary high-impact label is required, it is defined from the
training distribution only.  The primary severity score combines
reported affected population and fatalities after logarithmic
transformation:
\begin{equation}
  S_i
  =
  z_{\mathrm{train}}
  \left[
    \log_{10}
    \left(
      A_i+1
    \right)
  \right]
  +
  z_{\mathrm{train}}
  \left[
    \log_{10}
    \left(
      D_i+1
    \right)
  \right],
  \label{eq:app_severity_score}
\end{equation}
where $D_i$ is the reported number of fatalities and
$z_{\mathrm{train}}[\cdot]$ denotes standardisation using training-set
mean and standard deviation.  An event is classified as high impact if
\begin{equation}
  z_i
  =
  \mathbb{I}
  \left(
    S_i
    \geq
    Q_{0.90}^{\mathrm{train}}(S)
  \right).
  \label{eq:app_high_impact}
\end{equation}
Using the training distribution prevents the threshold from being
influenced by held-out test events.  Sensitivity analyses may also use
affected population alone, fatalities alone, or alternative quantile
thresholds.

\subsection{Missing data, uncertainty, and sensitivity checks}
\label{app:missing_uncertainty}

Missing continuous covariates are imputed using the median estimated
from the training set.  For each covariate $x_{i,j}$, a binary
missingness indicator is added:
\begin{equation}
  m_{i,j}
  =
  \mathbb{I}
  \left(
    x_{i,j}
    \ \mathrm{is\ missing}
  \right).
  \label{eq:app_missing_indicator}
\end{equation}
This allows the model to retain information about observational
availability while avoiding deletion of data-scarce African events.
No imputation parameter is estimated from the validation or test set.

Uncertainty in reported performance is quantified by event-level
bootstrap resampling.  For a metric $\mathcal{M}$ and bootstrap
replicates $b=1,\ldots,B$, the 95\% interval is
\begin{equation}
  \mathrm{CI}_{95}
  \left(
    \mathcal{M}
  \right)
  =
  \left[
    Q_{0.025}
    \left(
      \mathcal{M}^{(1)},\ldots,\mathcal{M}^{(B)}
    \right),
    Q_{0.975}
    \left(
      \mathcal{M}^{(1)},\ldots,\mathcal{M}^{(B)}
    \right)
  \right].
  \label{eq:app_bootstrap_ci}
\end{equation}
Sensitivity tests perturb the physics weight
$\lambda_{\mathrm{phys}}$, the wet-cell threshold $h_0$, Manning
roughness assumptions, reservoir size, and missing-data treatment.  A
result is considered robust only if the full physics-informed model
outperforms the rainfall-only and unconstrained baselines across the
main validation splits and remains stable under these perturbations.




% -----------------------------------------------------------------------
\clearpage
\section*{Appendix Tables and Figures}

\begin{table}[p]
\centering
\small
\caption{
African flood-study domains used for transfer evaluation.  Event
counts refer to the harmonised flood-event archive after quality
control.  Final values should be updated after the last event-screening
step.
}
\label{tab:app_study_domains}
\begin{tabular}{p{0.22\linewidth}p{0.24\linewidth}
                p{0.20\linewidth}p{0.22\linewidth}}
\toprule
Macro-region & Principal basin system & Approximate domain
& Hydrological rationale \\
\midrule
West Africa
&
Niger--Benue
&
$4^{\circ}$--$15^{\circ}$N,
$12^{\circ}$W--$15^{\circ}$E
&
Monsoon-driven flood seasonality, large floodplain storage,
rapidly increasing urban and peri-urban exposure, and recurrent
riverine flooding. \\

East Africa
&
Nile headwaters and Sudan--Ethiopia corridor
&
$4^{\circ}$S--$16^{\circ}$N,
$28^{\circ}$--$40^{\circ}$E
&
Strong seasonal rainfall contrasts, long routing pathways,
large basin memory effects, and high exposure along major river
corridors. \\

Southern Africa
&
Limpopo--Zambezi
&
$27^{\circ}$--$8^{\circ}$S,
$20^{\circ}$--$37^{\circ}$E
&
Cyclone-affected floods, reservoir-influenced river systems,
semi-arid drought--flood transitions, and strong interannual
hydroclimatic variability. \\
\botrule
\end{tabular}
\end{table}

\begin{table}[p]
\centering
\small
\caption{
Principal data groups used in the African flood experiments.
Observed flood footprints are reserved for validation and are not used
to construct exposure predictors.
}
\label{tab:app_data_sources}
\begin{tabular}{p{0.19\linewidth}p{0.27\linewidth}
                p{0.24\linewidth}p{0.20\linewidth}}
\toprule
Data group & Example products or sources & Variables extracted
& Role in the experiment \\
\midrule
Flood inventory and impacts
&
EM-DAT and complementary disaster-event inventories
\citep{EMDAT2024}
&
Event date, location, affected population, fatalities,
displacement, damages where available
&
Defines event--region records and impact targets. \\

Satellite flood observations
&
Global Flood Database and event-based optical or SAR flood maps
\citep{Tellman2021}
&
Observed inundation extent, wet/dry masks, flood duration where
available
&
Provides independent spatial validation targets only. \\

Hydroclimatic forcing
&
ERA5 or comparable ECMWF reanalysis products
\citep{Dee2011,Hersbach2020}
&
Daily precipitation, maximum precipitation, temperature,
soil moisture, evapotranspiration, runoff proxies
&
Defines meteorological forcing and antecedent hydroclimatic
memory. \\

Hydrological state
&
Gauge records where available, routed discharge products,
or hydrological reanalysis
\citep{VanDerKnijff2010,Alfieri2013,Thiemig2015}
&
Pre-event discharge anomaly, river-state percentile,
basin-storage indicator
&
Represents antecedent river condition and routing memory. \\

Terrain and landscape
&
SRTM, corrected DEM products, hydrographic networks,
land-cover maps
\citep{Farr2007,Yamazaki2017,Hawker2022,Lehner2013,ESA_CCI2017}
&
Elevation, slope, flow accumulation, drainage area, river
distance, land-cover class, roughness proxy
&
Defines routing controls, flood-prone supports, and spatial
heterogeneity. \\

Exposure and built environment
&
Gridded population, built-up layers, settlement products
\citep{Dobson2000,Stevens2015,Pesaresi2023}
&
Population count, built-up fraction, settlement density,
river-buffer population
&
Quantifies ex-ante human exposure before the event. \\
\botrule
\end{tabular}
\end{table}

\end{appendices}

\bibliography{flood-references}

\end{document}
